\documentclass[UTF8,12pt,a4paper,fontset=fandol]{ctexbook}


% 支持从项目根目录或当前笔记目录编译。
\makeatletter
\def\input@path{{../}}
\makeatother
\usepackage{researchnotes}


\title{凸优化笔记}
\author{}
\date{\today}


\begin{document}


\maketitle
\tableofcontents
\clearpage


\chapter*{课程大纲与阅读说明}
\addcontentsline{toc}{chapter}{课程大纲与阅读说明}

\section*{课程定位}

凸优化（convex optimization）研究在凸可行集上极小化凸函数的问题。本书面向学过线性代数与多元微积分的本科生，围绕三个问题展开：怎样建立凸模型，怎样证明候选解最优，以及怎样以有限精度求解。基本概念和核心算法给出推导；涉及较多凸分析工具的进阶定理注明引用或证明思路，供第二轮学习使用。

全书沿用 Boyd 和 Vandenberghe 的《Convex Optimization》\cite{boyd-vandenberghe} 的理论、应用、算法主线，以最小二乘、资源分配和几何投影贯穿各章。概率知识主要用于统计应用，数值线性代数按算法需要补充。

\section*{全书课程大纲}
\begin{enumerate}
  \item 第\ref{ch:co-introduction}章：绪论与优化建模。变量、目标、约束、可行性、最优值与数值精度。
  \item 第\ref{ch:co-sets}章：凸集。凸组合、凸包、典型凸集、保凸运算、投影、分离与对偶锥。
  \item 第\ref{ch:co-functions}章：凸函数。凸性判据、保凸运算、共轭、次梯度与拟凸性。
  \item 第\ref{ch:co-models}章：凸优化问题与标准模型。最优性、等价变换、线性规划、二次规划及锥规划。
  \item 第\ref{ch:co-duality}章：拉格朗日对偶与最优性条件。强弱对偶、Slater 条件、KKT 条件与灵敏度。
  \item 第\ref{ch:co-fitting}章：逼近、拟合与正则化。误差度量、稀疏拟合、鲁棒约束与形状限制。
  \item 第\ref{ch:co-statistics}章：统计估计与实验设计。似然、后验、熵、检测器和信息矩阵。
  \item 第\ref{ch:co-geometry}章：几何优化与分类。投影、距离、椭球、间隔与定位。
  \item 第\ref{ch:co-unconstrained}章：梯度方法与牛顿方法。线搜索、梯度法、投影与近端方法、次梯度法及牛顿法。
  \item 第\ref{ch:co-equality}章：等式约束优化算法。消元、约束牛顿法、残差与块线性系统。
  \item 第\ref{ch:co-interior}章：内点法。障碍函数、中心路径、可行性与原对偶迭代。
\end{enumerate}
附录\ref{app:co-background}补充数学基础，附录\ref{app:co-quadratic}讨论两个二次函数相关问题，附录\ref{app:co-numerical}说明数值线性代数。

\section*{阅读与练习}

首先顺序学习第\ref{ch:co-introduction}至\ref{ch:co-duality}章。应用部分的三个章节可以按兴趣选读；算法部分从梯度法及其处理简单约束、非光滑项的延伸出发，再学习无约束、等式约束和不等式约束牛顿方法。相对内部、矩阵凸性、自协调分析和附录\ref{app:co-quadratic}属于进阶内容，不影响首次掌握主要方法。

阅读时须分别回答：结论的假设是什么，关键不等式来自哪里，去掉某个假设是否仍然成立。各章习题给出必要提示或核对结果，完整推导仍应独立完成。数值练习要求同时检查目标值、可行性和适用的最优性指标；实验结果不能替代理论证明。

本书所说的“理解一个方法”，包括能用自己的话解释它的目标，能在一个两变量问题上手算，并能指出它何时不适用。例如，学完凸集应能判断两种可行方案的平均是否仍可行；学完凸函数应能解释为什么切线给出下界；学完对偶应能给一个候选解附上数值相等的下界；学完算法应能说明停止条件究竟控制什么误差。公式是这些判断的简明表达，而不是需要孤立记忆的规则。

首次学习可将每章的定义、基本定理、数值例题和习题作为主线。遇到一般分离定理、共轭双重化、自协调性等选读内容，先理解结论的用途及假设，再回到主线；不必为了尚未用到的推广中断学习。附录中的矩阵知识也可按需查阅。例如第一次看到 $P\succeq0$ 时，只需先记住它等价于所有方向上的二次型 $d^{\mathsf T}Pd$ 非负。

\chapter{绪论与优化建模}
\label{ch:co-introduction}

优化从一个选择问题开始：在允许的方案中，按给定准则选择最好的一种。数学模型必须分别说明哪些量可以改变、哪些限制必须满足，以及用什么量衡量优劣。

先看一个只有一个变量的选择。某次测量给出数值 $3$，希望选择与它接近的非负数 $x$，同时设备要求 $x\leq2$。用平方误差衡量接近程度，问题就是在区间 $[0,2]$ 上最小化 $(x-3)^2$。没有约束时答案为 $3$；加入上界后答案变为 $2$。这个例子已经包含优化的三个组成部分：$x$ 是可以选择的量，$(x-3)^2$ 是比较方案的准则，$[0,2]$ 是允许选择的范围。以后变量可以有几万个，但这三个问题始终不变。

\section{研究对象与基本语言}

设决策变量为 $x\in\mathbb R^n$，公共定义域为非空集合 $D\subseteq\mathbb R^n$。给定目标函数 $f_0:D\to\mathbb R$、不等式约束函数 $f_i:D\to\mathbb R$ 和等式约束函数 $h_j:D\to\mathbb R$，一般优化问题写为
\begin{equation}
 \begin{aligned}
 \mathop{\mathrm{minimize}}_{x\in D}\quad &f_0(x)\\
 \mathrm{subject\ to}\quad &f_i(x)\leq0,\quad i=1,\ldots,m,\\
 &h_j(x)=0,\quad j=1,\ldots,p.
 \end{aligned}
 \label{eq:co-general-problem}
\end{equation}
这里 $m,p$ 为约束数量，允许为零。$D$ 不可省略，例如含 $\log x_i$ 的函数要求 $x_i>0$。

\begin{definition}
满足全部约束且属于 $D$ 的点称为\keyterm{可行点（feasible point）}，其集合称为可行集 $\mathcal F$。最优值与最优解集分别为
\[
 p^\star=\inf_{x\in\mathcal F}f_0(x),\qquad
 X^\star=\{x\in\mathcal F:f_0(x)=p^\star\}.
\]
若 $x^\star\in X^\star$，则称它为全局最优解（global optimum）。若存在 $\delta>0$，使所有满足 $x\in\mathcal F$、$\|x-x^\star\|_2<\delta$ 的点都有 $f_0(x)\geq f_0(x^\star)$，则称 $x^\star$ 为局部最优解（local optimum）。
\end{definition}

约定 $\inf\varnothing=+\infty$。因此不可行问题有 $p^\star=+\infty$；若目标在可行集上无下界，则 $p^\star=-\infty$。有限最优值也可能不取到：$\inf_{x\in\mathbb R}e^x=0$，但不存在 $e^x=0$ 的实数。相反，常数目标在非空可行集上的每个点都是最优解。可行、有限、存在最优解和最优解唯一，是四个不同性质。

下确界可以理解为所有目标值共同承认的最大下界。例如在 $x>0$ 上最小化 $x$，所有可行目标值都大于零，而任意正数都不是下界，因为还可以取它的一半。因此下确界为零，却没有达到零的可行方案。这个区别会在统计估计和内点法中再次出现：一串方案越来越好，并不自动意味着它趋向某个可行最优点。

\begin{example}[四种问题状态]
比较以下问题：在 $x\geq1$ 且 $x\leq0$ 下最小化 $x^2$；在实轴上最小化 $x$；在 $x>0$ 上最小化 $x$；在 $0\leq x\leq1$ 上最小化常数零。
\end{example}
\begin{solve}
第一个问题没有可行点；第二个问题可行，但目标可以任意小；第三个问题有有限下确界而没有最优解；第四个问题中整个区间都是最优解集。因而“求解器没有给出唯一数值”可能对应完全不同的数学原因。分析时应先问可行性，再问有无下界、能否取到以及是否唯一。
\end{solve}

\section{从实际问题到数学模型}

建模通常依次确定决策变量与单位、目标、约束及数据。最大化收益可转为最小化其相反数；这种变换改变目标值的符号，但保留最优方案。

\begin{example}[资源分配]
两种产品的产量为 $x_1,x_2\geq0$，单位收益为 $3,2$，资源限制为 $x_1+x_2\leq4$、$2x_1+x_2\leq6$。求最大收益。
\end{example}
\begin{solve}
模型为最大化 $3x_1+2x_2$，满足上述四个线性不等式。对任意可行点，
\[
 3x_1+2x_2=(x_1+x_2)+(2x_1+x_2)\leq4+6=10.
\]
点 $(2,2)$ 可行且收益为 $10$，故它最优。两条资源约束的加权和给出了上界；这正是第\ref{ch:co-duality}章对偶证书的基本思想。
\end{solve}

硬约束要求条件必须满足；惩罚项允许违反条件并支付代价。例如约束 $x=0$ 与在目标中加入 $\alpha x^2$ 一般不等价，有限 $\alpha>0$ 未必迫使最优点为零。优化模型的正确求解也不保证建模假设符合实际。

以最小化 $(x-2)^2$ 为例，硬约束 $x=0$ 只允许答案零；若改为最小化 $(x-2)^2+\alpha x^2$，求导得到 $x_\alpha=2/(1+\alpha)$，对每个有限 $\alpha>0$ 都不为零。惩罚越大，违反原要求的程度越小，但这与严格满足约束仍有差别。另一方面，若把产量限制为整数，前面的资源分配就成为整数规划；只有当连续产量有实际意义，或者明确接受连续松弛时，才可直接使用凸模型。

变量的单位还决定系数的含义。如果一个变量以千克计、另一个以吨计，直接惩罚 $x_1^2+x_2^2$ 就隐含了不对称的偏好。归一化不是单纯的记号整理：在保留原问题的变量替换中，目标与约束的系数必须同时变换；若只改变正则项的尺度，则实际改变了模型。

\section{最小二乘与线性规划两个原型}

给定 $A\in\mathbb R^{r\times n}$、$b\in\mathbb R^r$，\keyterm{最小二乘（least squares）}以残差平方和衡量误差：
\begin{equation}
 \min_{x\in\mathbb R^n}\frac12\|Ax-b\|_2^2.
 \label{eq:co-least-squares}
\end{equation}
直接展开可得梯度 $A^{\mathsf T}(Ax-b)$，故候选解满足正规方程 $A^{\mathsf T}Ax=A^{\mathsf T}b$。若 $\widehat x$ 满足此式，则对任意 $d$，
\[
 \frac12\|A(\widehat x+d)-b\|_2^2
 =\frac12\|A\widehat x-b\|_2^2+\frac12\|Ad\|_2^2,
\]
交叉项因正规方程为零，因此 $\widehat x$ 最优。最优解总存在，因为 $b$ 到有限维闭子空间 $\operatorname{range}A$ 的正交投影存在；满列秩时 $A^{\mathsf T}A\succ0$，解唯一且为 $(A^{\mathsf T}A)^{-1}A^{\mathsf T}b$，秩亏时沿 $\ker A$ 移动不改变残差。

\keyterm{线性规划（linear programming，LP）}的目标和约束均为仿射函数，例如 $\min c^{\mathsf T}x$，满足 $Gx\leq h$、$Ax=b$。前一节资源分配即为 LP。最小二乘靠二次曲率刻画误差，LP 靠多面体约束限制方案，两者都属于凸优化。

\begin{example}[一条拟合直线的完整计算]
在位置 $t=0,1,2$ 观测到 $b=1,2,2$，用直线 $g(t)=u+vt$ 作最小二乘拟合。
\end{example}
\begin{solve}
参数是截距和斜率 $x=(u,v)^{\mathsf T}$，不是观测位置。设计矩阵、正规方程及解为
\[
 A=\begin{bmatrix}1&0\\1&1\\1&2\end{bmatrix},\quad
 b=\begin{bmatrix}1\\2\\2\end{bmatrix},\qquad
 \begin{bmatrix}3&3\\3&5\end{bmatrix}
 \begin{bmatrix}u\\v\end{bmatrix}
 =\begin{bmatrix}5\\6\end{bmatrix},\qquad
 u=\frac76,\quad v=\frac12.
\]
残差为 $(1/6,-1/3,1/6)^{\mathsf T}$，与 $A$ 的两列都正交，目标值为 $\|Ax-b\|_2^2/2=1/12$。任意其他直线都只能在这个正交残差上增加一个属于 $\operatorname{range}A$ 的向量；两部分平方长度相加，因此不能进一步减小误差。这给正规方程提供了几何解释。
\end{solve}

\section{凸优化的核心特征与边界}

集合是凸的，意指其中两点间的线段仍在集合内；函数是凸的，意指沿线段的函数值不超过端点函数值的加权平均。正式定义见第\ref{ch:co-sets}、\ref{ch:co-functions}章。凸问题的局部最优解必为全局最优解，但解的存在性和唯一性需要另外分析。

\keyterm{等价改写（equivalent reformulation）}应保持可行方案或最优方案的可恢复对应。\keyterm{凸松弛（convex relaxation）}则通常扩大可行集，用较易求解的问题提供界。例如对 $x\in\{0,1\}$ 的限制改用 $0\leq x\leq1$，极小化问题的最优值只能下降或不变，所得解未必满足原来的整数要求。

“凸”描述结构，计算成本还取决于规模、稀疏性、数据条件和精度。数值算法返回的通常是近似解：当 $p^\star$ 有限时，可行点 $x$ 若满足 $f_0(x)-p^\star\leq\varepsilon$，称为 $\varepsilon$-次优解。由于 $p^\star$ 未知，需要利用对偶下界等可计算量认证精度。

例如某个可行方案的成本为 $10.03$，另一个计算给出所有可行方案的成本都不小于 $10$，便得到最优值位于 $[10,10.03]$，当前方案距最优最多差 $0.03$。无需知道精确最优解，也可以作出有依据的停止判断。本书将反复使用这种“可行方案给上界，证书给下界”的思路。

\section{全书符号与表述约定}

向量默认为列向量，内积为 $x^{\mathsf T}y$，欧氏范数为 $\|x\|_2=\sqrt{x^{\mathsf T}x}$。$\|x\|_1=\sum_i|x_i|$，$\|x\|_\infty=\max_i|x_i|$。$x_k$ 表示迭代点，$(x_k)_i$ 表示其第 $i$ 个分量。$\mathbb R_+^n$、$\mathbb R_{++}^n$ 分别表示分量非负、严格为正的向量。

实对称矩阵空间为 $\mathbb S^n$。$X\succeq0$ 表示 $v^{\mathsf T}Xv\geq0$ 对所有 $v$ 成立；$X\succ0$ 表示该值对非零 $v$ 严格为正。相应集合为 $\mathbb S_+^n$、$\mathbb S_{++}^n$；$X\succeq Y$ 指 $X-Y\succeq0$，不同于逐元素比较。矩阵内积为 $\langle X,Y\rangle=\operatorname{tr}(X^{\mathsf T}Y)$，其中 $\operatorname{tr}$ 为迹。

函数的有效定义域记为 $\operatorname{dom}f$，梯度和海塞矩阵（Hessian matrix）记为 $\nabla f$、$\nabla^2f$。符号 $\inf$ 表示下确界，$\min$ 还断言该值取到；$\operatorname{arg\,min}$ 表示最优解集合。随机变量用大写字母，观测值通常用相应小写字母。

\section{例题与学习检查}
\begin{example}[均值作为最小二乘解]
给定实数观测 $b_1,\ldots,b_N$，求使 $\sum_{i=1}^N(x-b_i)^2$ 最小的常数 $x$。
\end{example}
\begin{solve}
令 $\bar b=N^{-1}\sum_i b_i$，利用 $\sum_i(b_i-\bar b)=0$ 得
\[
 \sum_i(x-b_i)^2=N(x-\bar b)^2+\sum_i(\bar b-b_i)^2.
\]
第二项与 $x$ 无关，第一项在且仅在 $x=\bar b$ 时为零，故均值是唯一最优解。
\end{solve}
\begin{enumerate}
 \item 分别构造不可行、无下界、有限但无最优解、存在多个最优解的问题。说明它们的 $p^\star$ 与 $X^\star$。
 \item 求 $\min_{x\geq1}(x-2)^2$ 和 $\min_{x\geq3}(x-2)^2$；比较最优点的导数。答案分别为 $2$ 和 $3$，约束边界最优点不必有零导数。
 \item 资源分配例中若第一项资源上限由 $4$ 改为 $4+u$，在 $u=0$ 附近求最优收益。提示：交点为 $(2-u,2+2u)$，收益为 $10+u$，须同时检查非负性。
 \item 基础练习：在 $0\leq x\leq2$ 上最小化 $(x-a)^2$，按 $a$ 所处区间给出答案。提示：最优解是把 $a$ 截断到 $[0,2]$，再解释它与硬约束的关系。
\end{enumerate}

\part{凸优化理论}

\chapter{凸集}
\label{ch:co-sets}

凸集描述可以连续混合的可行方案。本章建立集合运算、投影和分离的工具；这些工具随后用于函数凸性、最优性条件和对偶理论。

\section{仿射结构与凸结构}

假设两份生产计划分别为 $x=(2,0)$ 与 $y=(0,4)$。按一半时间执行各计划，平均产量为 $(1,2)$；按四分之一与四分之三混合，则为 $(1/2,3)$。表达式 $\theta x+(1-\theta)y$ 正是这种混合的数学形式。若所有比例的混合都仍满足原要求，就说允许的方案集合是凸的。这里混合的是方案本身；若方案包含不可分割的设备数量，混合得到的小数可能没有实际意义。

\begin{definition}
集合 $C\subseteq\mathbb R^n$ 为仿射集（affine set），若任意 $x,y\in C$、$\theta\in\mathbb R$ 都有 $\theta x+(1-\theta)y\in C$。若只要求
\begin{equation}
 x,y\in C,\quad0\leq\theta\leq1
 \quad\Longrightarrow\quad\theta x+(1-\theta)y\in C,
 \label{eq:co-convex-set}
\end{equation}
则称 $C$ 为凸集（convex set）。
\end{definition}

有限和 $\sum_{i=1}^k\theta_i x_i$ 在 $\sum_i\theta_i=1$ 时称为仿射组合；再有 $\theta_i\geq0$ 时称为凸组合。凸集包含其任意有限个点的凸组合：对点数归纳，将最后一个点与其余点的归一化凸组合再次作两点凸组合即可；权重为零的点可先删去。

集合 $S$ 的凸包（convex hull）$\operatorname{conv}S$ 是其所有有限凸组合组成的集合。两个凸组合的凸组合仍是凸组合，所以凸包为凸集；任何包含 $S$ 的凸集都包含这些组合，故凸包是包含 $S$ 的最小凸集。仿射包 $\operatorname{aff}S$ 同理由全部仿射组合构成。

仿射组合允许负权重，所以可以越过端点；凸组合只能位于端点之间。例如 $(1-t)(0,0)+t(1,0)$ 在 $t\in\mathbb R$ 时描出整条横轴，在 $0\leq t\leq1$ 时只描出一条线段。线性子空间必须经过原点，仿射集可以平移后不经过原点，而凸集还允许有边界和弯曲轮廓。这些概念描述的是逐步放宽的几何结构。

\begin{example}[从三个点得到一个三角形]
求 $S=\{(0,0),(2,0),(0,2)\}$ 的凸包，并写成不等式形式。
\end{example}
\begin{solve}
凸组合为 $(2\theta_2,2\theta_3)$，其中 $\theta_i\geq0$、$\theta_1+\theta_2+\theta_3=1$。令坐标为 $(u,v)$，得到 $u\geq0$、$v\geq0$、$u+v\leq2$。反过来，任何满足这些不等式的点都对应权重 $(1-(u+v)/2,u/2,v/2)$，所以两种表示完全等价。顶点表示便于理解“混合”，不等式表示便于检查一个方案是否可行。
\end{solve}

\begin{example}
仿射方程组的解集 $\{x:Ax=b\}$ 是仿射集；单位圆周不是凸集，因为两个相反点的中点为原点。其凸包为单位闭圆盘：盘内非零点可由同一直径的两个端点凸组合得到，原点取相等权重即可。
\end{example}

\section{凸锥与典型集合}

非空集合 $K$ 若对 $x\in K$、$t\geq0$ 有 $tx\in K$，称为锥（cone）；若还为凸集，则为凸锥（convex cone）。等价地，凸锥对非负线性组合封闭。以下集合反复出现：
\begin{itemize}
 \item 半空间 $\{x:a^{\mathsf T}x\leq b\}$ 与超平面 $\{x:a^{\mathsf T}x=b\}$，其中 $a\ne0$；有限个半空间和超平面的交称为多面体（polyhedron）。
 \item 范数球 $\{x:\|x-c\|\leq r\}$，其中 $r\geq0$。三角不等式保证两点的凸组合仍在球内。椭球可写为 $\{c+Bu:\|u\|_2\leq1\}$；$B$ 可逆时它为满维椭球。
 \item 概率单纯形 $\Delta_n=\{x\in\mathbb R_+^n:\mathbf1^{\mathsf T}x=1\}$，其中 $\mathbf1$ 为全一向量；二阶锥 $\mathcal Q_{n+1}=\{(x,t):\|x\|_2\leq t\}$。
 \item 半正定锥 $\mathbb S_+^n$：若 $X,Y\succeq0$，则对 $\alpha,\beta\geq0$ 和任意 $v$，有 $v^{\mathsf T}(\alpha X+\beta Y)v\geq0$，故它为凸锥。
\end{itemize}
多面体不一定有界。非负正交锥 $\mathbb R_+^n$ 是多面体凸锥；半正定矩阵的对角元非负是必要条件，但不是半正定性的充分条件。

在二维中，矩阵 $\left[\begin{smallmatrix}1&2\\2&1\end{smallmatrix}\right]$ 的对角元均为正，但沿 $v=(1,-1)^{\mathsf T}$ 的二次型为 $-2$，所以它不半正定。对称矩阵 $\left[\begin{smallmatrix}a&c\\c&b\end{smallmatrix}\right]$ 半正定的完整条件是 $a\geq0$、$b\geq0$、$ab-c^2\geq0$。矩阵序关心所有方向，而不只是坐标轴方向。

\section{保凸运算与集合表示}

\begin{proposition}
凸集的任意交、笛卡尔积、仿射像、仿射原像和集合和均为凸集。
\end{proposition}
\begin{proof}
若两点属于一族凸集的交，其线段属于族中每个集合，因而属于交。仿射映射 $T(x)=Ax+b$ 满足 $T(\theta x+(1-\theta)y)=\theta T(x)+(1-\theta)T(y)$，所以像和原像均保凸。对集合和 $C+D=\{c+d:c\in C,d\in D\}$，把两点凸组合分别在 $C,D$ 中展开即可；积集逐分量使用凸性。
\end{proof}

并集通常不保凸，例如 $\{-1\}\cup\{1\}$。投影是线性像，因此消去某些坐标保持凸性，但不一定保持闭性。

这里的坐标投影是“忘记部分坐标”，与后文“寻找距离最近的点”不同。例如集合
\[
 C=\{(u,v):u>0,\ v\geq1/u\}
\]
凸且闭：凸性来自 $1/u$ 的凸性；若其点列在平面内收敛，极限第一坐标不可能为零，否则第二坐标必须趋于无穷，因此极限仍满足约束。然而将第二坐标消去后得到 $(0,+\infty)$，它并不闭。这个例子说明辅助变量被消去后，凸性容易保留，最小值能否取到却仍要重新检查。

透视映射 $P(x,t)=x/t$ 定义于 $t>0$。若 $(x_1,t_1),(x_2,t_2)$ 属于此半空间内的凸集，为得到其像的权重 $\theta$，取原空间权重
\[
 \alpha=\frac{\theta t_2}{\theta t_2+(1-\theta)t_1}.
\]
代入可得 $P(\alpha(x_1,t_1)+(1-\alpha)(x_2,t_2))=\theta x_1/t_1+(1-\theta)x_2/t_2$，故透视像保凸。线性分式映射 $(Ax+b)/(c^{\mathsf T}x+d)$ 在分母为正的区域内是仿射映射与透视映射的复合，也保凸。

\section{相对内部、闭性与投影预备}

点 $x$ 为集合 $C$ 的内点，若某个以 $x$ 为中心的开球包含于 $C$。相对内部（relative interior）$\operatorname{ri}C$ 改为只要求开球与 $\operatorname{aff}C$ 的交包含于 $C$。例如 $\Delta_n$ 在 $n>1$ 时的内部为空，而相对内部由全部分量正的概率向量组成。集合包含其所有极限点时称为闭集；闭包 $\overline C$ 是包含 $C$ 的最小闭集。

具体地，$\Delta_2$ 是连接 $(1,0)$ 和 $(0,1)$ 的线段。在平面中，以其中任意一点为中心的圆盘都会离开直线 $x_1+x_2=1$，所以它没有平面内点；若只沿这条直线移动，中点附近的一小段都在集合里，因此中点属于相对内部。引入这个概念，是为了在含等式约束的问题中正确表达“还有活动余地”，而不是把低维可行集误判为完全没有内部结构。

闭性控制能否把极限点纳入集合，凸性控制两点之间能否连线。它们相互独立：开球凸但不闭，圆周闭但不凸。下面最近点定理同时使用闭性与凸性，前者帮助得到存在性，后者保证唯一性。

\begin{theorem}[最近点存在唯一性]
设 $C\subseteq\mathbb R^n$ 非空、闭且凸。对每个 $z\in\mathbb R^n$，存在唯一的 $u\in C$ 使 $\|z-u\|_2=\inf_{x\in C}\|z-x\|_2$。
\label{thm:co-projection-existence}
\end{theorem}
\begin{proof}
取 $x_0\in C$。寻找最近点时可限制于 $C\cap\{x:\|x-z\|_2\leq\|x_0-z\|_2\}$，这是非空紧集，连续距离函数在其上取得最小值。若不同的 $u,v$ 均为最近点，记共同平方距离为 $d^2$，则中点属于 $C$ 且
\[
 \left\|z-\frac{u+v}{2}\right\|_2^2
 =\frac12\|z-u\|_2^2+\frac12\|z-v\|_2^2-\frac14\|u-v\|_2^2<d^2,
\]
矛盾。
\end{proof}

\section{分离与支撑超平面}

平面中的直线 $a^{\mathsf T}x=\beta$ 把点按投影到法向量 $a$ 后的大小分成两侧，高维的超平面作用相同。分离定理要说明的是：如果一个点确实不属于闭凸集，就能找到一个线性不等式，集合内的点全部满足它，而这个点违反它。因而几何上的“在集合外”可以用一个线性证书表达。

\begin{theorem}[点与闭凸集的严格分离]
设 $C$ 非空闭凸且 $z\notin C$。存在 $a\ne0$、$\beta\in\mathbb R$，使 $a^{\mathsf T}x\leq\beta<a^{\mathsf T}z$ 对所有 $x\in C$ 成立。
\end{theorem}
\begin{proof}
取最近点 $u$。对任意 $x\in C$，点 $u+t(x-u)$ 在 $0\leq t\leq1$ 时属于 $C$，其到 $z$ 的平方距离在 $t=0$ 处有非负右导数，故 $(z-u)^{\mathsf T}(x-u)\leq0$。令 $a=z-u\ne0$、$\beta=a^{\mathsf T}u$，则 $a^{\mathsf T}z-\beta=\|z-u\|_2^2>0$。
\end{proof}

一般分离定理断言：两个不相交的非空凸集 $C,D$ 存在 $a\ne0$、$\beta$，使 $a^{\mathsf T}x\leq\beta\leq a^{\mathsf T}y$ 对所有 $x\in C,y\in D$ 成立。此有限维定理在这里引用，证明见\cite{boyd-vandenberghe}关于分离的讨论。若一个集合紧、另一个闭，则距离取到且为正，最近点对可给出严格分离。仅有闭性而无紧性，不足以保证正距离。

闭凸集 $C$ 的边界点 $u$ 存在支撑超平面（supporting hyperplane）：有 $a\ne0$ 使 $a^{\mathsf T}(x-u)\leq0$ 对所有 $x\in C$ 成立。可取趋于 $u$ 的集外点，将上述分离向量归一化后取收敛子列；极限向量仍为单位向量，极限不等式给出结论。

例如单位闭球外的点 $z=(2,0)$ 的最近点为 $u=(1,0)$，法向量 $z-u=(1,0)$ 给出支撑线 $x_1=1$。球内各点满足 $x_1\leq1$，而 $z$ 满足 $x_1=2$。分离定理的证明把这个直观构造推广到了任意非空闭凸集，并不要求边界光滑。

\section{广义不等式与对偶锥}

正常锥（proper cone）是闭、凸、有非空内部且满足 $K\cap(-K)=\{0\}$ 的锥。定义 $x\preceq_K y$ 当且仅当 $y-x\in K$，则它具有自反性、传递性和反对称性，但两点可能不可比较。向量逐分量序和对称矩阵半正定序都是例子。集合的最小元须小于等于每个元素；极小元只要求不存在另一个不等于自身且小于等于自身的元素。

凸锥 $K$ 的对偶锥（dual cone）定义为
\begin{equation}
 K^*=\{y:y^{\mathsf T}x\geq0\ \text{对所有 }x\in K\}.
 \label{eq:co-dual-cone}
\end{equation}
它是闭半空间的交，故为闭凸锥。$\mathbb R_+^n$ 自对偶，因为可取各坐标单位向量检验符号。二阶锥也自对偶：若 $\|y\|_2\leq s$、$\|x\|_2\leq t$，则 $y^{\mathsf T}x+st\geq0$；反之取边界向量检验可得 $s\geq\|y\|_2$。$\mathbb S_+^n$ 在迹内积下自对偶：对 $X,Y\succeq0$，$\operatorname{tr}(XY)=\operatorname{tr}(X^{1/2}YX^{1/2})\geq0$；反向取 $X=vv^{\mathsf T}$ 即得 $v^{\mathsf T}Yv\geq0$。

\begin{proposition}
闭凸锥满足 $K^{**}=K$；一般凸锥满足 $K^{**}=\overline K$。
\end{proposition}
\begin{proof}
由定义 $K\subseteq K^{**}$。若 $z\notin K$，严格分离给出 $a^{\mathsf T}x\leq\beta<a^{\mathsf T}z$。锥的缩放性质迫使 $a^{\mathsf T}x\leq0$ 对所有 $x\in K$，且因 $0\in K$ 有 $\beta\geq0$。于是 $-a\in K^*$ 而 $(-a)^{\mathsf T}z<0$，所以 $z\notin K^{**}$。非闭情形利用 $K^*=(\overline K)^*$ 即得。
\end{proof}

\section{例题与学习检查}
\begin{example}[到半空间的投影]
设 $a\ne0$，$C=\{x:a^{\mathsf T}x\leq b\}$。若 $a^{\mathsf T}z\leq b$，最近点就是 $z$；否则最近点为
\[
 u=z-\frac{a^{\mathsf T}z-b}{\|a\|_2^2}a.
\]
事实上 $a^{\mathsf T}u=b$，且任意可行 $x$ 满足 $(z-u)^{\mathsf T}(x-u)\leq0$。展开 $\|z-x\|_2^2$ 即知它不小于 $\|z-u\|_2^2$。
\end{example}
\begin{enumerate}
 \item 判断 $\{x:\|Ax-b\|_2\leq1\}$、$\{x:\|x\|_2=1\}$ 是否凸，并指出退化维数可能带来的例外。
 \item 求 $\{(x,y):x\geq0,\ 0\leq y\leq x\}$ 的对偶锥。提示：原锥由 $(1,0),(1,1)$ 生成，对偶条件为 $u\geq0,u+v\geq0$。
 \item 证明闭凸集的相交可以没有内点，即使每个集合都有内点；再计算所得集合的相对内部。
 \item 综合练习：将三角形 $\operatorname{conv}\{(0,0),(2,0),(0,2)\}$ 与半空间 $x_1+x_2\geq1$ 相交，写出顶点与不等式两种表示。核对结果为四边形，其顶点是 $(1,0),(2,0),(0,2),(0,1)$。
\end{enumerate}

\chapter{凸函数}
\label{ch:co-functions}

凸函数把集合的线段性质转化为函数值的不等式。判断凸性既可以求导，也可以把复杂函数分解为已知凸函数及保凸运算；后者往往更适合建模。

\section{定义、上图与次水平集}

对平方函数，两个输入的平均所对应的输出，不超过两个输出的平均。例如 $f(x)=x^2$ 时，$f((0+2)/2)=1$，而 $(f(0)+f(2))/2=2$。几何上，曲线在两点之间位于弦的下方；在多元情形，则沿每一条线段检查同样的性质。“凸”并不等于单调递增，$x^2$ 在负半轴递减而在正半轴递增，却在整个实轴上都是凸函数。

\begin{definition}
设 $D\subseteq\mathbb R^n$ 为凸集。函数 $f:D\to\mathbb R$ 为凸函数（convex function），若
\begin{equation}
 f(\theta x+(1-\theta)y)\leq\theta f(x)+(1-\theta)f(y),
 \quad x,y\in D,\quad\theta\in[0,1].
 \label{eq:co-convex-function}
\end{equation}
若对 $x\ne y$、$0<\theta<1$ 不等式严格，则称为严格凸函数（strictly convex function）。若 $-f$ 凸，则 $f$ 为凹函数（concave function）。
\end{definition}

上图（epigraph，亦称上境图）与水平 $\alpha$ 下的次水平集分别为
\[
 \operatorname{epi}f=\{(x,t):x\in D,\ f(x)\leq t\},\qquad
 C_\alpha=\{x\in D:f(x)\leq\alpha\}.
\]
$f$ 凸当且仅当其上图凸：正向对两个上图点应用式\eqref{eq:co-convex-function}，反向取高度恰为 $f(x),f(y)$ 的两点即可。凸函数的次水平集也凸，因为两点函数值均不超过 $\alpha$ 时，其凸组合的函数值仍不超过 $\alpha$。

允许函数取 $+\infty$，可将 $D$ 外的点统一排除。对于 $f:\mathbb R^n\to\mathbb R\cup\{+\infty\}$，有效定义域为 $\operatorname{dom}f=\{x:f(x)<+\infty\}$；定义域非空时称为适当函数（proper function）。上图闭时称为闭函数（closed function），在有限维空间中等价于下半连续。集合 $C$ 的指示函数（indicator function）$I_C$ 在 $C$ 上为零、在其外为 $+\infty$；它为凸函数当且仅当 $C$ 凸。

例如在 $x\geq0$ 上最小化 $(x+1)^2$，可以写成在全空间最小化 $(x+1)^2+I_{[0,+\infty)}(x)$。负数的代价为无穷，因此不会成为有限目标的候选解。这里的指示函数与概率论中取值为零或一的示性函数不同。扩展值记号统一了目标与约束，但不能因此忽略定义域：即使某个表达式局部二阶导数非负，若指定的定义域不凸，也不能直接称它为本书意义下的凸函数。

\section{一阶、二阶判据与曲率}

\begin{theorem}[微分判据]
设 $D$ 为开凸集。若 $f$ 在 $D$ 上可微，则它凸当且仅当
\begin{equation}
 f(y)\geq f(x)+\nabla f(x)^{\mathsf T}(y-x),\quad x,y\in D.
 \label{eq:co-tangent-bound}
\end{equation}
若 $f$ 二次连续可微，则它凸当且仅当 $\nabla^2f(x)\succeq0$ 对所有 $x\in D$ 成立。
\end{theorem}
\begin{proof}
由凸性，$[f(x+t(y-x))-f(x)]/t\leq f(y)-f(x)$。令 $t\downarrow0$ 得式\eqref{eq:co-tangent-bound}。反向令 $z=(1-\theta)x+\theta y$，分别在 $z$ 处对 $x,y$ 使用切平面下界，以 $1-\theta,\theta$ 加权后梯度项抵消，即得凸性。

二阶判据可沿直线证明。令 $g(t)=f(x+td)$，则 $g''(t)=d^{\mathsf T}\nabla^2f(x+td)d$。海塞矩阵半正定使每条线上的导数单调，积分即得一阶下界；反之凸的二次可微单变量函数有 $g''(0)\geq0$，对任意方向 $d$ 即得到矩阵半正定。
\end{proof}

若海塞矩阵处处正定，则函数严格凸；逆命题不成立，$x^4$ 严格凸但在零点的二阶导数为零。给定 $\mu>0$，若 $f(x)-\mu\|x\|_2^2/2$ 凸，则称 $f$ 为 $\mu$-强凸函数（strongly convex function）。可微情形等价于
\[
 f(y)\geq f(x)+\nabla f(x)^{\mathsf T}(y-x)+\frac\mu2\|y-x\|_2^2.
\]
在二次连续可微情形，强凸性等价于 $\nabla^2f(x)\succeq\mu I$。另一方面，梯度的 $L$-利普希茨连续性（Lipschitz continuity）指 $\|\nabla f(x)-\nabla f(y)\|_2\leq L\|x-y\|_2$，它控制梯度变化的上界。对开凸域上的二次连续可微凸函数，这等价于 $\nabla^2f(x)\preceq LI$。

一阶判据的意义比“二阶导数非负”更直接：在任意点作出的切平面，都给出全局下界。以 $f(x)=x^2$ 为例，在 $x=1$ 处的下界是 $1+2(y-1)=2y-1$，两者之差为 $(y-1)^2$。强凸性进一步要求这个差至少包含一个统一的正二次项，因此它给出了离开最优点后目标至少增加多少的定量信息。

\begin{example}[三个容易混淆的曲率性质]
比较实轴上的 $x^2$、$x^4$ 和 $e^x$ 的严格凸性、强凸性与全局梯度利普希茨性。
\end{example}
\begin{solve}
$x^2$ 的二阶导数恒为 $2$，所以它是 $2$-强凸函数，梯度以 $L=2$ 利普希茨连续。$x^4$ 的导数 $4x^3$ 严格递增，所以函数严格凸；但二阶导数在零点为零，不存在正的全局强凸常数，又因 $12x^2$ 无上界，不存在全局梯度利普希茨常数。$e^x$ 的二阶导数处处为正，所以严格凸；它在负无穷方向趋于零、在正无穷方向无界，因此也没有这两个全局常数。若限制到有界区间，部分结论会改变，故曲率界必须连同适用区域一起陈述。
\end{solve}

检查多元函数的海塞矩阵时，只看各个二阶偏导 $\partial^2 f/\partial x_i^2$ 不够。例如 $f(x,y)=x^2+y^2+4xy$ 的两个对角二阶偏导均为 $2$，但沿 $(1,-1)$ 的二阶方向导数为 $-4$，函数并不凸。混合偏导描述变量之间的耦合，不能省略。

\section{常用函数与詹森不等式}

仿射函数同时凸、凹。二次函数 $f(x)=\tfrac12x^{\mathsf T}Px+q^{\mathsf T}x+r$ 在 $P=P^{\mathsf T}$ 时凸当且仅当 $P\succeq0$。任意范数都凸，因为三角不等式和正齐次性给出定义中的不等式。$e^x$ 在实轴上凸，$-\log x$ 在 $x>0$ 上凸，二者可直接用二阶导数验证。

对数和指数函数（log-sum-exp）$f(x)=\log\sum_{i=1}^n e^{x_i}$ 的梯度分量为 $q_i=e^{x_i}/\sum_j e^{x_j}$，海塞矩阵为 $\operatorname{diag}(q)-qq^{\mathsf T}$。由于 $q_i\geq0$、$\sum_iq_i=1$，对任意 $v$，
\[
 v^{\mathsf T}\nabla^2f(x)v
 =\sum_iq_i v_i^2-\left(\sum_iq_i v_i\right)^2
 =\sum_iq_i\left(v_i-\sum_jq_jv_j\right)^2\geq0,
\]
故 $f$ 凸。它满足 $\max_i x_i\leq f(x)\leq\max_i x_i+\log n$，可用作最大值的光滑近似。

在 $X\succ0$ 上，$F(X)=-\log\det X$ 沿对称方向 $H$ 的二阶导数为
\[
 D^2F(X)[H,H]=\operatorname{tr}(X^{-1}HX^{-1}H)
 =\|X^{-1/2}HX^{-1/2}\|_{\mathrm F}^2\geq0,
\]
其中 $\|B\|_{\mathrm F}^2=\operatorname{tr}(B^{\mathsf T}B)$ 为弗罗贝尼乌斯范数平方，所以 $F$ 凸。

\begin{proposition}[詹森不等式]
若 $f$ 凸，$\theta_i\geq0$、$\sum_i\theta_i=1$，则
\[
 f\left(\sum_i\theta_i x_i\right)\leq\sum_i\theta_i f(x_i).
\]
\end{proposition}
\begin{proof}
对点数归纳，将最后一点与其余点的归一化凸组合作两次应用即可。若 $f$ 严格凸，等号要求所有正权重点相同，否则在归纳中至少一次出现严格不等式。
\end{proof}
期望形式称为 Jensen 不等式。一个便于核查的版本是：$f$ 在开凸集 $D$ 上凸且可微，随机向量 $Z$ 几乎处处属于 $D$，$\mathbb EZ\in D$，且 $Z,f(Z)$ 可积，则 $f(\mathbb EZ)\leq\mathbb E f(Z)$。在 $\mathbb EZ$ 处使用式\eqref{eq:co-tangent-bound}并取期望，线性项的期望为零，即得结论。

例如对正数 $a,b$，将凹函数 $\log$ 的 Jensen 不等式应用于相等权重，得到 $\log((a+b)/2)\geq(\log a+\log b)/2$，指数化后就是算术平均不小于几何平均。对具有有限二阶矩的随机变量，取 $f(z)=z^2$，则 $(\mathbb EZ)^2\leq\mathbb EZ^2$，二者之差正是方差。凸性不等式由此把几何、代数与概率联系起来。

\section{保凸运算与复合规则}

非负加权和与仿射复合保凸。若每个 $f_\alpha$ 凸，则 $f(x)=\sup_\alpha f_\alpha(x)$ 为凸扩展值函数，因为
\[
 f_\alpha(\theta x+(1-\theta)y)
 \leq\theta\sup_\beta f_\beta(x)+(1-\theta)\sup_\beta f_\beta(y)
\]
对每个 $\alpha$ 成立。逐点最小值不保凸，例如 $\min\{x,-x\}=-|x|$。

若外层函数 $h$ 凸且单调不减、内层 $g$ 凸，则在复合有定义且定义域凸的条件下，$h\circ g$ 凸：先用 $g$ 的凸性及 $h$ 的单调性，再用 $h$ 的凸性。若 $h$ 单调不增，则应配合凹的 $g$。多变量外层函数按各分量的单调方向检验，不能只检查外层凸性。

设 $F(x,y)$ 联合凸，$C$ 为固定凸集，$f(x)=\inf_{y\in C}F(x,y)$，且下确界不为 $-\infty$。取使 $F(x_i,y_i)\leq f(x_i)+\varepsilon$ 的近似最优点，使用联合凸性，再令 $\varepsilon\downarrow0$，便得 $f$ 凸。是否取到下确界不是此证明的必要条件，但 $f$ 未必闭。

凸函数 $f$ 的透视函数（perspective）为 $g(x,t)=t f(x/t)$，定义于 $t>0$、$x/t\in\operatorname{dom}f$。令 $t=\theta t_1+(1-\theta)t_2$，则 $x/t$ 是 $x_1/t_1,x_2/t_2$ 以 $\theta t_1/t,(1-\theta)t_2/t$ 为权的凸组合，代入凸性并乘以 $t$ 即证 $g$ 联合凸。因此 $\|x\|_2^2/t$ 在 $t>0$ 上联合凸。

\begin{example}[判断复合函数时先检查单调性]
判断 $e^{x^2}$、$(x^2-1)^2$ 以及 $\sqrt{1+x^2}$ 在实轴上的凸性。
\end{example}
\begin{solve}
指数函数凸且递增，因此 $e^{x^2}$ 可直接用复合规则判为凸。平方函数虽然凸，却在整个实轴上不单调，故不能仅凭“平方套平方”判断第二个函数；实际二阶导数为 $12x^2-4$，在零点附近为负，所以不凸。第三个函数的外层平方根是凹的，常用复合规则不能直接给出凸性；但它等于二维向量 $(1,x)$ 的欧氏范数，是范数与仿射映射的复合，故凸。规则不能证明，和函数不凸，是两回事。
\end{solve}

\section{共轭函数、次梯度与支撑下界}

先从一个固定斜率 $y$ 的仿射下界 $y^{\mathsf T}x-c$ 出发。要使它对所有 $x$ 都不超过 $f(x)$，必须有 $c\geq y^{\mathsf T}x-f(x)$ 对所有 $x$ 成立。因此给定斜率时，最紧的这种下界所需的常数，就是下面定义的 $f^*(y)$。共轭函数把“按位置描述函数”转为“按斜率描述其支撑下界”。

函数 $f$ 的凸共轭（convex conjugate）为
\begin{equation}
 f^*(y)=\sup_{x\in\operatorname{dom}f}\{y^{\mathsf T}x-f(x)\}.
 \label{eq:co-conjugate}
\end{equation}
它是仿射函数的上确界，故总是凸的。例如 $f(x)=\|x\|_2^2/2$ 时，配方得 $f^*(y)=\|y\|_2^2/2$。$I_C^*(y)=\sup_{x\in C}y^{\mathsf T}x$ 称为 $C$ 的支撑函数（support function）。由定义立即得到 Fenchel 不等式 $f(x)+f^*(y)\geq y^{\mathsf T}x$。

\begin{definition}
设 $f$ 为适当凸函数，$x\in\operatorname{dom}f$。若向量 $g$ 满足 $f(z)\geq f(x)+g^{\mathsf T}(z-x)$ 对所有 $z$ 成立，则称 $g$ 为 $f$ 在 $x$ 处的次梯度（subgradient）；全部次梯度组成次微分（subdifferential）$\partial f(x)$。
\end{definition}

当 $x$ 位于有效定义域的内部且凸函数 $f$ 在 $x$ 可微时，次微分为 $\{\nabla f(x)\}$；$f(x)=|x|$ 的次微分在正点为 $\{1\}$、负点为 $\{-1\}$、零点为 $[-1,1]$。这由 $|z|\geq gz$ 对所有 $z$ 的条件直接得到。Fenchel 不等式取等号当且仅当 $y\in\partial f(x)$，因为等号正好要求 $x$ 在共轭定义的上确界中达到最大值。

次梯度保留的是“支撑下界”而不一定是某个方向上的普通导数。在 $|x|$ 的尖点零处，斜率 $-1$ 到 $1$ 之间的每条过原点直线都在图像下方，因此次梯度是一个区间。对 $f(x)=\max_i(a_i^{\mathsf T}x+b_i)$，若第 $i$ 项在 $x$ 处达到最大值，则对任意 $z$，
\[
 f(z)\geq a_i^{\mathsf T}z+b_i
 =f(x)+a_i^{\mathsf T}(z-x),
\]
故 $a_i$ 是一个次梯度；多个活跃斜率的凸组合也都是次梯度。这使分段线性函数也能用一阶信息分析。

\begin{example}[绝对值的共轭]
求实轴上 $f(x)=|x|$ 的共轭函数。
\end{example}
\begin{solve}
若 $|y|\leq1$，则 $yx-|x|\leq0$ 对所有 $x$ 成立，取 $x=0$ 达到零。若 $y>1$，令 $x\to+\infty$；若 $y<-1$，令 $x\to-\infty$，均使表达式无界增大。因此 $f^*(y)=0$ 当 $|y|\leq1$，其余处为 $+\infty$，即区间 $[-1,1]$ 的指示函数。一个惩罚项在对偶中可以变成一个约束，这是正则化对偶模型的重要来源。
\end{solve}

次微分在定义域边界上可能为空。例如 $f(x)=-\sqrt{x}$ 定义于 $x\geq0$ 并在域外取 $+\infty$，在零点若有有限次梯度 $g$，则 $-1/\sqrt z\geq g$ 对所有 $z>0$，与 $z\downarrow0$ 矛盾。

作为选读结论，有限维 Fenchel--Moreau 定理断言：适当闭凸函数满足 $f^{**}=f$。它说明该函数可由全部仿射下界恢复；一般函数只能直接得到 $f^{**}\leq f$。此定理可由上图分离证明，本节仅引用结论；该结果与闭性条件见\cite{boyd-vandenberghe}的共轭函数小节。

\section{拟凸性、对数凸性与矩阵凸性}

函数 $f$ 称为拟凸（quasiconvex），若每个次水平集都凸；等价地，$f(\theta x+(1-\theta)y)\leq\max\{f(x),f(y)\}$。凸函数必拟凸，反之不成立：$x^3$ 在实轴上的次水平集均为半直线，但其二阶导数在负半轴为负。拟凸最小化可通过逐次判断水平集是否与约束集相交来求解。

正值函数 $f$ 若 $\log f$ 凸，称为对数凸（log-convex）；若 $\log f$ 凹，称为对数凹（log-concave）。例如 $e^{x^2}$ 对数凸，$e^{-x^2}$ 对数凹。对数凸函数是凸函数，因为 $f=\exp(\log f)$ 是凸递增函数与凸函数的复合。

对矩阵值函数 $F:D\to\mathbb S^r$，关于半正定序的凸性指
\[
 F(\theta x+(1-\theta)y)\preceq\theta F(x)+(1-\theta)F(y).
\]
它等价于对每个 $v\in\mathbb R^r$，标量函数 $v^{\mathsf T}F(x)v$ 凸。例如 $F(x)=xx^{\mathsf T}$ 为此意义下的凸映射，因为右端减左端为 $\theta(1-\theta)(x-y)(x-y)^{\mathsf T}\succeq0$。逐元素凸性与矩阵序凸性不能混同。

\section{例题与学习检查}
\begin{example}[曲率与正则化]
对 $\alpha>0$，函数 $f(x)=\tfrac12\|Ax-b\|_2^2+\tfrac\alpha2\|x\|_2^2$ 的海塞矩阵为 $A^{\mathsf T}A+\alpha I\succeq\alpha I$，因此它至少为 $\alpha$-强凸，即使 $A$ 秩亏也如此。若将平方正则项换成 $\alpha\|x\|_1$，目标仍凸，但一般不再可微或强凸。
\end{example}
\begin{enumerate}
 \item 用海塞矩阵和透视函数两种方法证明 $x^2/t$ 在 $t>0$ 上联合凸。
 \item 求 $f(x)=\tfrac12x^{\mathsf T}Px$ 在 $P\succ0$ 时的共轭。答案为 $\tfrac12y^{\mathsf T}P^{-1}y$。
 \item 证明 $\max_i(a_i^{\mathsf T}x+b_i)$ 凸；若当前最大值由唯一指标取得，求其梯度。
 \item 举例说明严格凸不推出强凸，凸不推出对数凸。提示：$x^4$ 与正半轴上的 $x$。
 \item 基础练习：求 $f(x)=\max\{2x,-x\}$ 在零点的次微分。答案为 $[-1,2]$；应分别检验 $z>0$ 与 $z<0$ 的支撑不等式。
 \item 推导练习：证明 $f(x)=\log(e^x+e^{-x})$ 严格凸，但在整个实轴上没有正的强凸常数。提示：$f''(x)=4/(e^x+e^{-x})^2$。
\end{enumerate}

\chapter{凸优化问题与标准模型}
\label{ch:co-models}

模型的凸性应从变量、定义域、目标及每个约束逐项检查。本章把前两章的判据转化为建模工具，并给出算法和对偶理论所需的基本最优性结论。

\section{标准形式与最优性基本定理}

凸目标需要和正确方向的约束配合。约束 $f_i(x)\leq0$ 选取凸函数的次水平集，因而保凸；约束 $f_i(x)\geq0$ 则没有同样保证。例如 $x^2\leq1$ 给出区间 $[-1,1]$，而 $x^2\geq1$ 给出两段分离的半直线。对等式约束，要求仿射是一个统一且容易检查的充分结构：它既不会像 $x^2=1$ 那样形成分离的点，也允许直接消元。

式\eqref{eq:co-general-problem}在 $D$ 凸、$f_0,f_1,\ldots,f_m$ 凸、$h_j$ 仿射时称为标准形式凸问题。它的可行集是凸集，因为凸函数次水平集与仿射集的交仍凸。最大化凹函数也属于凸优化，因为取负号后得到凸的极小化问题。

一个凸可行集可能被写成非标准约束。例如 $x^2=0$ 定义单点凸集，但左侧不是仿射函数；改写为 $x=0$ 就得到标准凸表示。建模时应寻找便于验证和求解的表示。

\begin{theorem}[凸问题的最优解]
在凸集 $C$ 上极小化凸函数 $f$ 时，每个局部最优解都是全局最优解；最优解集是凸集；若 $f$ 严格凸，则最优解至多一个。
\end{theorem}
\begin{proof}
若局部最优点 $x^\star$ 不是全局最优点，存在 $y\in C$ 使 $f(y)<f(x^\star)$。对任意 $0<t<1$，$x_t=(1-t)x^\star+ty$ 可行且 $f(x_t)<f(x^\star)$；令 $t$ 足够小，便与局部最优性矛盾。若 $x,y$ 均最优，其凸组合的函数值不超过最优值，又不可能小于它，故仍最优。严格凸时两个不同最优点的中点会有更小函数值，故不能同时存在。
\end{proof}

存在性需另外保证。连续函数在非空紧可行集上取得最小值；更一般地，下半连续目标只要有一个非空紧的有限次水平集，就能在其上取得全局最小值。若目标在闭可行集上连续且随 $\|x\|_2\to\infty$ 趋于 $+\infty$，这种紧次水平集也存在。严格凸性只保证至多一个解，不能代替这些条件。

局部即全局的证明依靠两个独立条件：可行集凸，使通向更好方案的短线段仍然可行；目标凸，使这条线段上立即出现目标改善。缺少前者，即使目标为 $(x-2)^2$，在可行集 $\{-1,2\}$ 上 $-1$ 也是局部最优却不是全局最优，因为它附近根本没有其他可行点。缺少后者，沿通向远处更好方案的线段，也可能先上升后下降。

\section{一阶最优性与约束边界}

\begin{theorem}
设 $C$ 非空凸，$f$ 在包含 $C$ 的开凸集上凸且可微。点 $x^\star\in C$ 最优，当且仅当
\begin{equation}
 \nabla f(x^\star)^{\mathsf T}(x-x^\star)\geq0,
 \quad\text{对所有 }x\in C.
 \label{eq:co-first-order-optimality}
\end{equation}
\end{theorem}
\begin{proof}
必要性由线段 $x^\star+t(x-x^\star)$ 上目标在 $t=0$ 处的右导数非负得到。充分性由切平面下界 $f(x)\geq f(x^\star)+\nabla f(x^\star)^{\mathsf T}(x-x^\star)\geq f(x^\star)$ 得到。
\end{proof}

当 $C=\mathbb R^n$ 时，可取 $x-x^\star=-\nabla f(x^\star)$，故条件化为 $\nabla f(x^\star)=0$。约束问题则不必如此：在 $x\geq0$ 上最小化 $f(x)=x$，最优点为零而导数为 $1$。可行方向限制了实际能够下降的方向。

对适当凸函数，包括通过指示函数编码约束的情形，由次梯度定义可知
\[
 x^\star\text{ 最优}\quad\Longleftrightarrow\quad0\in\partial f(x^\star).
\]
把它进一步拆成目标和约束各自的次微分时，需要相应的求和规则及资格条件。

一阶最优性还给出一个可以计算的目标下界。对当前可行点 $x$，若
\[
 G_C(x)=\sup_{y\in C}\nabla f(x)^{\mathsf T}(x-y)
\]
有限，则 $0\leq f(x)-p^\star\leq G_C(x)$：将切平面下界对 $y\in C$ 取下确界即可。这个量比较的是当前点与所有可行方向的线性改善空间。在有界多面体上计算它是一个线性规划；在无界集合上它可能无穷大，因此并非任何问题都适用。

\section{等价变换、松弛与拟凸优化}

上图改写 $\min_x f(x)$ 为 $\min_{x,t}t$，约束 $f(x)\leq t$，保留最优值，并可由新解直接恢复 $x$。对于 $\min_x\max_i f_i(x)$，改写为 $\min t$、$f_i(x)\leq t$，常能消除目标中的非光滑最大值。可逆变量替换保留方案对应；严格递增的目标变换保留最优解，但通常改变最优值。

辅助变量不会自动带来等价性，必须检查两个方向。上图改写中，任意原可行 $x$ 都能配上 $t=f(x)$，而任意新可行 $(x,t)$ 又满足 $f(x)\leq t$，因此双方的最优值相同。若写成 $f(x)=t$，虽然也能保持最优值，但非仿射凸函数的等式图像通常不是凸集；写成不等式正是保留凸结构的关键。

\begin{example}[极小极大问题的手算与线性规划]
选择一个实数 $x$，使它到 $0$ 与 $4$ 的最坏距离尽量小，即最小化 $\max\{|x|,|x-4|\}$。
\end{example}
\begin{solve}
引入 $t$ 后，约束为 $-t\leq x\leq t$、$-t\leq x-4\leq t$。其中 $x\leq t$ 和 $x\geq4-t$ 要同时成立，必须有 $t\geq2$。取 $x=2,t=2$ 满足全部约束，所以最优值为 $2$。既可以把答案理解为两个端点的中点，也可以把 $t\geq2$ 理解为原问题的最优性证书。
\end{solve}

若以更大的集合 $\widetilde C\supseteq C$ 替代原可行集，则
\[
 \inf_{x\in\widetilde C}f(x)\leq\inf_{x\in C}f(x).
\]
因此松弛产生下界；只有松弛解还满足原约束时，才能直接认证它也是原问题的最优解。

设 $f$ 拟凸、约束集 $C$ 凸，且能判断 $C\cap\{x:f(x)\leq\tau\}$ 是否为空。给定最优值下界 $l$ 和来自可行点的上界 $u$，每次取 $\tau=(l+u)/2$：可行时保存该点并令 $u=\tau$，不可行时令 $l=\tau$。经过 $\lceil\log_2((u_0-l_0)/\varepsilon)\rceil$ 次判定，区间长度不超过 $\varepsilon$；此计数不包括每次可行性判定的代价，并假设判定精确。

\section{线性规划与凸二次规划}

LP 的常用标准形式为
\[
 \min_x c^{\mathsf T}x,\qquad Ax=b,\quad x\geq0.
\]
自由变量可写成两个非负变量之差；不等式 $Gx\leq h$ 可引入非负松弛变量 $s$，写成 $Gx+s=h$。转换增加变量，但便于统一分析。

凸二次规划（quadratic programming，QP）具有形式
\[
 \min_x\ \tfrac12x^{\mathsf T}Px+q^{\mathsf T}x+r,
 \qquad Gx\leq h,\quad Ax=b,\quad P\succeq0.
\]
若不等式也为凸二次函数，则称为凸二次约束二次规划（quadratically constrained quadratic programming，QCQP）。每一个不等式的二次型矩阵都应半正定；任意二次等式通常不符合标准凸形式。最小二乘为 QP，因为其二次项矩阵为 $A^{\mathsf T}A$。

\section{二阶锥规划与半正定规划}

二阶锥规划（second-order cone programming，SOCP）的目标线性，约束由仿射等式及
\[
 \|A_i x+b_i\|_2\leq c_i^{\mathsf T}x+d_i
\]
组成。左、右两侧整体表示一个仿射向量落入二阶锥；右侧非负性已由该不等式保证。凸 QCQP 可以通过二次型分解和旋转二阶锥转为 SOCP，例如 $\|u\|_2^2\leq t$ 等价于 $\|(2u,t-1)\|_2\leq t+1$。

半正定规划（semidefinite programming，SDP）允许仿射矩阵约束 $F_0+\sum_i x_iF_i\succeq0$，其中 $F_i$ 为给定对称矩阵。最大特征值极小化可写为 $\min t$，满足 $tI-F(x)\succeq0$，前提是 $F$ 关于变量仿射。

\begin{proposition}[Schur 补]
若 $A\succ0$，则
\[
 \begin{bmatrix}A&B\\B^{\mathsf T}&C\end{bmatrix}\succeq0
 \quad\Longleftrightarrow\quad C-B^{\mathsf T}A^{-1}B\succeq0.
\]
\end{proposition}
\begin{proof}
对向量 $u,v$，块二次型等于
\[
 (u+A^{-1}Bv)^{\mathsf T}A(u+A^{-1}Bv)
 +v^{\mathsf T}(C-B^{\mathsf T}A^{-1}B)v.
\]
充分性由两项非负得到；必要性取 $u=-A^{-1}Bv$。
\end{proof}
若 $A$ 仅半正定，须增加值域条件并使用广义逆，不能直接照用此公式。

SOCP 与 SDP 的名称来自所使用的凸锥，而不是目标函数外观的复杂程度。例如 $x_1^2+x_2^2\leq t$ 表面上含有平方项，仍可用一个二阶锥约束表示。验证等价式时，先由 $\|(2u,t-1)\|_2\leq t+1$ 知道右端非负，再平方得 $4\|u\|_2^2+(t-1)^2\leq(t+1)^2$，即 $\|u\|_2^2\leq t$。反向中 $t\geq\|u\|_2^2\geq0$ 保证平方不会引入符号错误。

一个最小的 SDP 例子是
\[
 \begin{bmatrix}t&x\\x&1\end{bmatrix}\succeq0
 \quad\Longleftrightarrow\quad t\geq x^2.
\]
因此矩阵不等式既能描述多个方向的曲率，也能编码熟悉的标量二次不等式。遇到矩阵模型时，可以先在 $2\times2$ 情形展开检查，再理解其高维推广。

\section{几何规划与对数变换}

对正变量 $x_j>0$，单项式（monomial）为 $c\prod_jx_j^{a_j}$，其中 $c>0$、$a_j\in\mathbb R$；若干单项式的和称为正项式（posynomial）。几何规划（geometric programming，GP）的标准形式为最小化正项式，满足正项式不等式 $f_i(x)\leq1$ 和单项式等式 $h_j(x)=1$。

令 $y_j=\log x_j$。单项式的对数成为 $\log c+a^{\mathsf T}y$，正项式的对数成为
\[
 \log\sum_k\exp(a_k^{\mathsf T}y+\log c_k),
\]
这是凸函数。因此对目标取对数、对每个约束取对数后，得到标准凸问题，原变量由 $x_j=e^{y_j}$ 恢复。

\begin{example}
长方形要求 $x,y>0$、面积 $xy\geq1$，求最小周长。模型为最小化 $2x+2y$，满足 $x^{-1}y^{-1}\leq1$，是 GP。由 $x+y\geq2\sqrt{xy}\geq2$，最优解为 $x=y=1$；对数改写中，面积约束成为 $-u-v\leq0$。
\end{example}

\section{多目标优化与建模检查}

设可行集为 $C$，目标向量为 $F(x)=(f_1(x),\ldots,f_q(x))$。若 $x^\star\in C$，且不存在 $x\in C$ 使各目标均不大于 $F(x^\star)$ 且至少一项严格更小，则称 $x^\star$ 为帕累托最优（Pareto optimal）。给定权重 $w_i>0$，求解 $\min_{x\in C}\sum_iw_i f_i(x)$，称为加权和标量化（scalarization）。其最优解必为帕累托最优，否则支配它的点会有严格更小的加权和。若权重非负且不全为零，一般只能保证弱帕累托最优，即不存在同时严格改善全部目标的可行点。

例如在 $0\leq x\leq1$ 上同时最小化 $x^2$ 和 $(x-1)^2$，每一点均为帕累托最优；正权重和的解为 $x=w_2/(w_1+w_2)$，展示了偏好权重对方案的影响。

一个模型的最终检查应包括定义域、目标曲率、约束方向、等式仿射性以及辅助变量能否还原。诸如 $\|x\|_2\geq1$ 的约束虽使用凸函数，却是其超水平集，通常非凸。

\section{例题与学习检查}
\begin{example}[同一残差的不同模型]
令 $e=Ax-b\in\mathbb R^r$。最小化 $\|e\|_1$ 等价于 LP：最小化 $\mathbf1^{\mathsf T}t$，约束 $-t\leq Ax-b\leq t$；最小化 $\|e\|_\infty$ 等价于最小化标量 $s$，约束 $-s\mathbf1\leq Ax-b\leq s\mathbf1$。最小化 $\|e\|_2$ 可用 SOCP 上图，而最小化其平方可用 QP。三种误差度量一般给出不同最优方案。
\end{example}
\begin{enumerate}
 \item 将 $\min_x\max\{|x-1|,2|x+1|\}$ 写成 LP 并求解。答案为 $x=-1/3$，最优值 $4/3$。
 \item 设 $f:\mathbb R^n\to\mathbb R$ 为 $\mu$-强凸连续函数，$C$ 为非空闭凸集。证明 $f$ 在 $C$ 上的最小值存在且唯一。提示：强凸下界提供强制性。
 \item 用 Schur 补将 $y^{\mathsf T}X^{-1}y\leq t$、$X\succ0$ 写成矩阵不等式，并说明正定定义域不可遗漏。
 \item 建模练习：把 $\min_x |x-1|+2|x+2|$ 写成 LP。再用次梯度核验 $x^\star=-2$、最优值为 $3$，比较“写成标准模型”与“解出模型”两个步骤。
\end{enumerate}

\chapter{拉格朗日对偶与最优性条件}
\label{ch:co-duality}

可行点给出极小化问题的上界；对偶理论构造下界。当上下界一致时，最优性得到认证。对偶变量还描述约束扰动的影响，并把几何最优性条件写成可以求解的方程与不等式。

回顾资源分配例，收益满足 $3x_1+2x_2\leq10$，并不是因为逐个列出了所有可行方案，而是因为把两条资源约束各乘以一再相加。对偶理论系统地寻找这种“加权后的约束证书”。对于极小化问题，方向改为构造下界；非负权重保证不等式相加时方向正确，等式的权重则可以取任意符号。

\section{拉格朗日函数与对偶函数}

对式\eqref{eq:co-general-problem}，引入不等式乘子 $\lambda\in\mathbb R_+^m$、等式乘子 $\nu\in\mathbb R^p$，定义拉格朗日函数（Lagrangian）和对偶函数（dual function）：
\begin{equation}
 \mathcal L(x,\lambda,\nu)=f_0(x)+\sum_i\lambda_i f_i(x)+\sum_j\nu_j h_j(x),
 \qquad g(\lambda,\nu)=\inf_{x\in D}\mathcal L(x,\lambda,\nu).
 \label{eq:co-lagrangian-dual-function}
\end{equation}
对 $x$ 取下确界时保留定义域 $D$，但不再要求已纳入拉格朗日函数的约束。$g$ 可以取 $-\infty$，此时该乘子不给出有用的有限下界。

为何要去掉这些约束再取下确界？因为在原可行点上，$\lambda_i f_i(x)\leq0$，所以拉格朗日函数不大于原目标；若再允许 $x$ 在更大的定义域内选择，下确界只会更小。这两次放松共同产生一个可靠的下界。直接把原可行约束全部保留在内层，虽然仍能定义一个下界，却通常没有简化求解，也不是这里所用的对偶构造。

\begin{example}[一个可以画出全部下界的对偶]
求问题 $\min x^2/2$、$x\geq2$ 的对偶函数，并比较几组乘子的作用。
\end{example}
\begin{solve}
将约束写为 $2-x\leq0$。对 $\lambda\geq0$，配方得到
\[
 \mathcal L(x,\lambda)=\frac12(x-\lambda)^2
                  +2\lambda-\frac12\lambda^2,\qquad
 g(\lambda)=2\lambda-\frac12\lambda^2.
\]
乘子 $0,1,2,3$ 分别给出下界 $0,3/2,2,3/2$。选一个合法乘子即可得到下界，对偶优化则在这些下界中选最大的一个。凹抛物线 $g$ 在 $\lambda=2$ 达到 $2$，而原可行点 $x=2$ 的目标也为 $2$，于是无需遍历可行半直线就证明了最优性。
\end{solve}

\begin{proposition}
对偶函数关于乘子为凹函数，无须假设原问题凸。
\end{proposition}
\begin{proof}
固定 $x$ 时，$\mathcal L$ 关于乘子仿射。用 $y=(\lambda,\nu)$ 简记乘子，对 $0\leq\theta\leq1$ 有
\[
 \inf_x\mathcal L(x,\theta y_1+(1-\theta)y_2)
 \geq\theta\inf_x\mathcal L(x,y_1)+(1-\theta)\inf_x\mathcal L(x,y_2),
\]
这正是凹性。
\end{proof}

\begin{example}
对 $\min_x x$，约束 $x\geq1$，有 $\mathcal L=x+\lambda(1-x)$。若 $\lambda\ne1$，对自由变量 $x$ 的线性函数无下界；若 $\lambda=1$，$g(1)=1$。对偶函数的有限性本身产生了乘子约束，不能只形式化地求导后忽略其余情况。
\end{example}

\section{对偶问题、弱对偶与证书}

对偶问题为最大化 $g(\lambda,\nu)$，满足 $\lambda\geq0$；其最优值为 $d^\star$。这是在凸域上最大化凹函数，因此是凸优化问题。

\begin{theorem}[弱对偶]
对任意原可行点 $x$ 与 $\lambda\geq0$、$\nu\in\mathbb R^p$，有 $g(\lambda,\nu)\leq f_0(x)$，从而 $d^\star\leq p^\star$。
\end{theorem}
\begin{proof}
由下确界定义及可行性，
\[
 g(\lambda,\nu)\leq\mathcal L(x,\lambda,\nu)
 =f_0(x)+\sum_i\lambda_i f_i(x)\leq f_0(x).
\]
分别对原变量取下确界、对乘子取上确界即可。
\end{proof}

当原、对偶值有限时，对偶间隙（duality gap）$f_0(x)-g(\lambda,\nu)$ 给出
\[
 0\leq f_0(x)-p^\star\leq f_0(x)-g(\lambda,\nu).
\]
因此原可行点与对偶可行点的目标差不超过 $\varepsilon$，就认证了 $\varepsilon$-次优性。这个结论不需要强对偶，但在最优对偶间隙为正时可能无法得到任意小的证书。

弱对偶还可以识别某些问题状态。如果对偶目标可以任意大，那么原问题不可能存在具有有限目标值的可行点；反过来，原目标无下界时，不可能存在一个给出有限下界的对偶点。这些结论都是同一不等式的推论。但“没有找到有限对偶点”本身不是原问题无下界的证明，计算失败与数学不可行必须区别开来。

\section{强对偶与 Slater 条件}

若 $d^\star=p^\star$，称为强对偶（strong duality）。凸性之外还需适当的约束资格条件（constraint qualification）来保证它。下面给出适合首次学习的版本。

强对偶说的是两侧的最优值相同，不是说任意一对原、对偶点的值都相同，也不是说两侧一定存在最优解。Slater 条件的几何含义是：在满足等式后，还能找到一点使每条不等式都有严格余量。这个余量防止可行集仅靠某种退化接触形成，从而帮助把分离超平面转化为有限的拉格朗日乘子。首次学习可以先掌握定理的条件和后面的反例，再阅读分离证明。

\begin{theorem}[Slater 条件的有限值版本]
设 $f_0,\ldots,f_m$ 在 $\mathbb R^n$ 上为有限凸函数，等式约束为 $Ax=b$，$A$ 满行秩。若存在 $\bar x$ 满足 $A\bar x=b$、$f_i(\bar x)<0$ 对所有 $i$ 成立，且 $p^\star$ 有限，则强对偶成立，且存在最优乘子。此结论不要求原最优值取到。
\label{thm:co-slater}
\end{theorem}
\begin{proof}
考虑凸集
\[
 \mathcal U=\{(u,v,t):\text{存在 }x,\ f_i(x)<u_i,\ Ax-b=v,\ f_0(x)<t\}.
\]
点 $(0,0,p^\star)$ 不属于 $\mathcal U$。由分离定理，存在非零向量 $(\lambda,\nu,\eta)$，使
\[
 \lambda^{\mathsf T}u+\nu^{\mathsf T}v+\eta t\geq\eta p^\star
 \quad\text{对所有 }(u,v,t)\in\mathcal U.
\]
$u_i,t$ 可任意增大，故 $\lambda\geq0,\eta\geq0$。若 $\eta=0$，在严格可行点 $\bar x$ 取 $f_i(\bar x)<u_i<0$、$v=0$，得 $\lambda=0$。随后 $\nu^{\mathsf T}(Ax-b)\geq0$ 对所有 $x$，迫使 $A^{\mathsf T}\nu=0$；满行秩推出 $\nu=0$，与分离向量非零矛盾。因此 $\eta>0$，归一化使其为 $1$。

固定任意 $x$，令 $u_i\downarrow f_i(x)$、$t\downarrow f_0(x)$，得到 $\mathcal L(x,\lambda,\nu)\geq p^\star$。取下确界知 $g(\lambda,\nu)\geq p^\star$；弱对偶给出反向不等式，故这组乘子取到 $d^\star=p^\star$。
\end{proof}

更一般的版本允许凸的有效定义域 $D$，要求严格可行点位于 $\operatorname{ri}D$；仿射不等式还可放宽为非严格成立。该推广见\cite{boyd-vandenberghe}关于 Slater 条件的讨论。本书使用时明确采用哪个版本。

资格条件是充分条件而非必要条件。例如 $\min_x x$、$x^2\leq0$ 的原最优值为零，对偶函数在 $\lambda>0$ 时为 $-1/(4\lambda)$，故强对偶仍成立，但有限乘子不能取到零。这里没有严格可行点，KKT 乘子也不存在。

作为对照，将目标改为 $x^2$ 而保留 $x^2\leq0$，仍然没有严格可行点，但取 $\lambda=0$ 时对偶值已为零，且原点满足 KKT 条件。因而 Slater 条件不成立，只表示这个充分定理不能使用，不能立即推断强对偶失败或乘子不存在。检查具体模型时，还可以尝试把 $x^2\leq0$ 改写为等价的仿射等式 $x=0$，消除由表示带来的退化。

\section{KKT 条件与互补松弛}

设 $D$ 为开凸集，各 $f_i$ 在其上凸且可微，等式约束为 $Ax=b$。Karush--Kuhn--Tucker 条件（KKT conditions）为
\begin{equation}
 \begin{aligned}
 &x^\star\in D,\quad f_i(x^\star)\leq0,\quad Ax^\star=b,\\
 &\lambda_i^\star\geq0,\quad\lambda_i^\star f_i(x^\star)=0,\quad i=1,\ldots,m,\\
 &\nabla f_0(x^\star)+\sum_i\lambda_i^\star\nabla f_i(x^\star)
    +A^{\mathsf T}\nu^\star=0.
 \end{aligned}
 \label{eq:co-kkt}
\end{equation}
前三类关系分别表达原可行性、对偶可行性和互补松弛（complementary slackness），最后一式为拉格朗日函数的驻点条件。

驻点式可以理解为一种力的平衡：目标的下降方向被若干活跃约束的法向量阻挡，等式乘子则补偿只能沿可行平面运动的限制。互补松弛把这种解释写成代数关系：若某条不等式还有余量，它就不需要提供阻挡作用，乘子必为零；若乘子为正，约束必须恰好卡在边界。活跃约束（active constraint）只表示 $f_i(x^\star)=0$，并不自动表示它的乘子严格为正。

\begin{theorem}
满足式\eqref{eq:co-kkt}的点是原、对偶最优点。反之，若原最优解存在、强对偶成立且对偶最优值取到，则原、对偶最优点满足 KKT 条件。因此在 Slater 条件下，KKT 是已有原最优解的必要充分刻画。
\end{theorem}
\begin{proof}
拉格朗日函数关于 $x$ 凸，驻点条件使 $x^\star$ 取到其下确界；互补松弛和原可行性给出 $g(\lambda^\star,\nu^\star)=f_0(x^\star)$，由弱对偶知最优。

反向，从 $g(\lambda^\star,\nu^\star)\leq\mathcal L(x^\star,\lambda^\star,\nu^\star)\leq f_0(x^\star)$ 且两端相等可知，中间两步均取等号。各 $\lambda_i^\star f_i(x^\star)\leq0$ 且总和为零，故逐项为零；$x^\star$ 又是开域内拉格朗日函数的极小点，所以梯度为零。
\end{proof}

若某约束严格不活跃，即 $f_i(x^\star)<0$，其乘子必为零；活跃约束的乘子却也可能为零。若把非空闭凸约束 $x\in C$ 保留在定义域中，则在 $x\in C$ 处定义法锥（normal cone）$N_C(x)=\{v:v^{\mathsf T}(z-x)\leq0\text{ 对所有 }z\in C\}$，相应的拉格朗日最小化条件为 $-\nabla_x\mathcal L\in N_C(x)$。乘子的存在仍须单独核查资格条件。

\begin{example}[先猜活跃约束，再用 KKT 核查]
\label{ex:co-active-constraints}
求
\[
 \min_{x,y}\ \frac12\bigl[(x-2)^2+(y-1)^2\bigr],
 \qquad x+y\leq2,\quad x\geq0,\quad y\geq0.
\]
\end{example}
\begin{solve}
无约束解 $(2,1)$ 违反总量限制，所以先猜总量约束活跃，而非负约束有余量。分别给三个不等式配乘子 $\lambda,\mu,\eta\geq0$，驻点式为
\[
 x-2+\lambda-\mu=0,\qquad
 y-1+\lambda-\eta=0.
\]
猜测 $x,y>0$ 后，互补松弛给 $\mu=\eta=0$；再用 $x+y=2$，解得 $\lambda=1/2$、$x=3/2$、$y=1/2$。结果确实非负、可行，并满足所有互补式，所以猜测得到验证，KKT 的充分性证明它全局最优，目标值为 $1/4$。目标严格凸，故解唯一。若算出的某个坐标为负，必须撤回对应的不活跃假设，重新讨论边界情形。

在点 $(1/2,1/2)$ 三条不等式都严格成立，因此 Slater 条件也成立，必要性保证没有遗漏不满足 KKT 的最优点。小规模问题可以这样枚举活跃情况；约束很多时，逐个枚举会产生过多组合，应使用后面的数值方法。
\end{solve}

KKT 中的驻点条件单独使用很容易出错。例如在 $x\geq1$ 上最小化 $x^2$，直接令目标导数为零会得到不可行点零；而在实轴上最小化非凸函数 $x^4-x^2$，驻点零反而是局部极大点。本节的全局充分性同时依靠凸性、原可行性、乘子符号、互补松弛和驻点条件，不能只保留最后一行。

\section{鞍点、共轭与不同对偶表示}

拉格朗日鞍点（saddle point）满足
\[
 \mathcal L(x^\star,\lambda,\nu)
 \leq\mathcal L(x^\star,\lambda^\star,\nu^\star)
 \leq\mathcal L(x,\lambda^\star,\nu^\star)
\]
对所有 $x\in D,\lambda\geq0,\nu$ 成立。右侧表示 $x^\star$ 最小化拉格朗日函数；左侧迫使原可行性和互补松弛，否则相应乘子可使左侧无界增大。因此鞍点与达到相同有限值的原对偶最优点对应。

对 $\min_x f(x)$、$Ax=b$，共轭给出
\[
 g(\nu)=\inf_x\{f(x)+(A^{\mathsf T}\nu)^{\mathsf T}x-b^{\mathsf T}\nu\}
       =-f^*(-A^{\mathsf T}\nu)-b^{\mathsf T}\nu.
\]
例如 $f(x)=\|x\|_2^2/2$ 时，$g(\nu)=-\|A^{\mathsf T}\nu\|_2^2/2-b^{\mathsf T}\nu$。若 $A$ 满行秩，则
\[
 \nu^\star=-(AA^{\mathsf T})^{-1}b,\qquad
 x^\star=A^{\mathsf T}(AA^{\mathsf T})^{-1}b.
\]
这就是等式约束最小范数解，原对偶值同为 $b^{\mathsf T}(AA^{\mathsf T})^{-1}b/2$。

LP $\min c^{\mathsf T}x$、$Ax=b$、$x\geq0$ 的对偶也可由有限性求得。将 $x\geq0$ 留在定义域中，则 $g(\nu)=-b^{\mathsf T}\nu$ 当且仅当 $c+A^{\mathsf T}\nu\geq0$，否则为 $-\infty$。改变变量 $y=-\nu$ 得熟悉的形式：最大化 $b^{\mathsf T}y$，满足 $A^{\mathsf T}y\leq c$。

\section{扰动、灵敏度与乘子的解释}

把约束改为 $f_i(x)\leq u_i$、$Ax=b+v$，最优值记为 $p(u,v)$。设原问题强对偶且有最优乘子 $(\lambda^\star,\nu^\star)$，则对扰动问题的任意可行点使用原拉格朗日下界，有
\begin{equation}
 p(u,v)\geq p(0,0)-{\lambda^\star}^{\mathsf T}u-{\nu^\star}^{\mathsf T}v.
 \label{eq:co-sensitivity-bound}
\end{equation}
确切地说，$f_0(x)\geq g(\lambda^\star,\nu^\star)-{\lambda^\star}^{\mathsf T}u-{\nu^\star}^{\mathsf T}v$，再对扰动可行集取下确界即可。若 $p$ 在原点邻域有限且可微，支撑下界的斜率唯一，因而 $\nabla_up(0,0)=-\lambda^\star$、$\nabla_vp(0,0)=-\nu^\star$。

乘子称为影子价格（shadow price），表示放宽单位约束可能带来的边际目标改进。等式扰动的符号取决于采用 $Ax=b+v$ 还是其他约定，不能脱离定义记忆。

在例\ref{ex:co-active-constraints}中，将总量上限改为 $2+u$。只要 $-1<u<1$，原来的活跃结构保持不变，解为
\[
 x(u)=\frac{3+u}{2},\qquad y(u)=\frac{1+u}{2},\qquad
 p(u)=\frac{(1-u)^2}{4}.
\]
因此 $p'(0)=-1/2=-\lambda^\star$，确实与乘子解释一致。这里精确值满足 $p(u)=1/4-u/2+u^2/4$，而对偶提供的是全局线性下界 $p(u)\geq1/4-u/2$。只有取足够小的扰动时，才可把线性项解释为近似变化量；扰动很大导致活跃约束改变时，不能继续把原来的乘子当作新的精确斜率。

\section{择一定理与锥对偶}

\begin{theorem}[Farkas 引理]
对 $A\in\mathbb R^{r\times n}$、$b\in\mathbb R^r$，下列两个系统恰有一个可行：
\[
 Ax=b,\ x\geq0;\qquad A^{\mathsf T}y\geq0,\ b^{\mathsf T}y<0.
\]
\end{theorem}
\begin{proof}
两者不能同时可行，否则 $b^{\mathsf T}y=x^{\mathsf T}A^{\mathsf T}y\geq0$。若第一个不可行，则 $b$ 不属于有限生成锥 $K=\{Ax:x\geq0\}$。该锥是闭的：每个锥组合可消去线性相关的活动生成元，写成线性无关生成元的非负组合；这样的子锥各自闭，且只有有限多个。由点与闭凸锥的严格分离，存在 $y\in K^*$ 满足 $b^{\mathsf T}y<0$，而 $y\in K^*$ 等价于 $A^{\mathsf T}y\geq0$。
\end{proof}

对于仿射锥约束 $F(x)\in-K$，选择 $z\in K^*$，则 $\langle z,F(x)\rangle\leq0$，故它可以作为拉格朗日项。对 SDP 约束 $F(x)\succeq0$，相应项写为 $-\langle Z,F(x)\rangle$，$Z\succeq0$；互补性为 $\langle Z,F(x)\rangle=0$。两者半正定时还可推出 $ZF(x)=0$，但不是逐元素相乘。

\section{例题与学习检查}
\begin{example}[用 KKT 求解与认证]
求 $\min_{x\geq1}x^2/2$。
\end{example}
\begin{solve}
$\mathcal L=x^2/2+\lambda(1-x)$，$g(\lambda)=\lambda-\lambda^2/2$，$\lambda\geq0$。KKT 给出 $x=\lambda$、$\lambda(1-x)=0$、$x\geq1$，因此 $x^\star=\lambda^\star=1$。原对偶目标均为 $1/2$，最优性得到认证。点 $x=2$ 严格可行，因此 Slater 条件也成立。
\end{solve}

\begin{example}[比例资源分配]
给定权重 $w_i>0$、预算 $B>0$，最大化 $\sum_iw_i\log x_i$，满足 $x_i>0$、$\sum_i x_i\leq B$。
\end{example}
\begin{solve}
转为最小化负效用。驻点条件为 $-w_i/x_i+\lambda=0$；由于效用对每个变量递增，预算约束取等号。记 $W=\sum_iw_i$，得到 $\lambda^\star=W/B$、$x_i^\star=Bw_i/W$。这些点满足 KKT，因此最优；负效用严格凸，故解唯一。预算的边际效用为 $W/B$，与乘子一致。
\end{solve}
\begin{enumerate}
 \item 求 $\min_{x\geq0}(x-a)^2/2$ 的最优解和乘子，按 $a$ 的符号讨论。答案为 $x^\star=\max\{a,0\}$、$\lambda^\star=\max\{-a,0\}$。
 \item 验证 $\min x$、$x^2\leq0$ 的最优点不满足任何有限乘子的 KKT 条件，解释它与充分性定理并不矛盾。
 \item 对任意原可行点和对偶可行点证明：若目标相等，它们分别最优；说明该论证为何不需要 Slater 条件。
 \item 综合练习：求 $\min_{x+y\leq B,\ x,y\geq0}[(x-2)^2+(y-1)^2]/2$ 对所有 $B>0$ 的解。提示：在 $0<B\leq1$ 时为 $(B,0)$；在 $1<B<3$ 时为 $((B+1)/2,(B-1)/2)$；在 $B\geq3$ 时为 $(2,1)$。分别给出合法乘子，观察活跃约束怎样改变。
\end{enumerate}

\part{凸优化应用}

\chapter{逼近、拟合与正则化}
\label{ch:co-fitting}

拟合问题需要回答两件事：什么样的误差不可接受，以及什么样的解更符合任务需要。前者决定损失函数，后者表现为约束或正则项。本章固定设计矩阵 $A\in\mathbb R^{r\times n}$ 与观测 $b\in\mathbb R^r$，以 $e(x)=Ax-b$ 记残差。

\section{范数逼近与误差度量}

最小二乘最小化 $\sum_i e_i^2$，对大残差施加较重惩罚；最小绝对偏差最小化 $\sum_i|e_i|$，对残差大小的惩罚为线性；极小极大逼近最小化 $\max_i|e_i|$，控制最坏单点误差。后两者的 LP 表示已在第\ref{ch:co-models}章给出。

最小化 $\|e\|_2$ 与 $\|e\|_2^2$ 有相同最优解，因为平方在非负实数上严格递增。不同范数一般没有这种对应。权重模型 $\sum_iw_i e_i^2$ 在 $w_i\geq0$ 时凸；权重可以表达量纲归一化或已知的观测可靠性，但应与建模假设相符。

平方损失对应的单点导数是 $2e_i$，残差越大，对更新方向的影响越大；绝对值损失的次梯度大小至多为一，因此单个大残差的影响受到限制。极小极大准则更关心目前最大的那几个误差，允许其他较小误差适当增大。选择损失之前，应先说清楚任务究竟关心平均表现、抗异常能力，还是最坏情况，不能仅凭某个模型容易计算来决定。

\begin{example}[均值、中位数与中程数]
用常数 $x$ 拟合观测 $0,0,9$。平方损失的唯一解为均值 $3$；绝对偏差损失为 $2|x|+|x-9|$，在 $x=0$ 唯一最小；最大绝对误差要求同时接近最小值 $0$ 和最大值 $9$，故最优解为 $9/2$，误差为 $9/2$。同一数据下，不同目标表达了不同偏好。
\end{example}

均值、中位数和中程数的差别也可以一般地推出。对有序数据 $b_{(1)}\leq\cdots\leq b_{(N)}$，绝对偏差和在不等于观测值的位置的导数为“左侧观测个数减右侧观测个数”。导数从负变正的位置就是中位数；若 $N$ 为偶数，两个中间观测之间的整个区间都最优。对于最大绝对偏差，只有两端 $b_{(1)},b_{(N)}$ 决定最坏距离，二者的中点给出最小半径 $(b_{(N)}-b_{(1)})/2$。这些结论说明非唯一性有时来自损失的平坦段，并不是算法没有收敛。

\section{最小范数解与带约束逼近}

欠定方程 $Ax=b$ 可能有很多解。最小范数模型 $\min\|x\|_2^2/2$、$Ax=b$ 选出唯一最小欧氏范数解；当 $A$ 满行秩时，第\ref{ch:co-duality}章已得
\[
 x^\star=A^{\mathsf T}(AA^{\mathsf T})^{-1}b.
\]
任意其他解为 $x^\star+d$，$Ad=0$；因 $x^\star\in\operatorname{range}A^{\mathsf T}$，有 ${x^\star}^{\mathsf T}d=0$，故 $\|x^\star+d\|_2^2=\|x^\star\|_2^2+\|d\|_2^2$，再次证明唯一性。一般可行的秩亏系统可用伪逆 $A^\dagger b$ 表示，见附录\ref{app:co-numerical}。

实际拟合常要求 $x\geq0$、$\mathbf1^{\mathsf T}x=1$ 或 $l\leq x\leq u$。这些约束仿射，故与凸损失组合后仍凸。例如非负最小二乘的 KKT 条件为 $A^{\mathsf T}(Ax-b)-\lambda=0$、$x,\lambda\geq0$、$x_i\lambda_i=0$。将无约束解的负分量直接截断通常不是其解，因为矩阵 $A^{\mathsf T}A$ 会耦合各分量。

\begin{example}[为什么不能总是先拟合再截断]
令 $A=\left[\begin{smallmatrix}1&1\\0&1\end{smallmatrix}\right]$、$b=(0,1)^{\mathsf T}$，求非负最小二乘解。
\end{example}
\begin{solve}
无约束方程 $Ax=b$ 的解为 $(-1,1)$，逐分量截断给 $(0,1)$，其目标为 $1/2$。在非负域上，固定 $x_2\geq0$ 时，第一项 $(x_1+x_2)^2/2$ 在 $x_1=0$ 最小；于是问题化为最小化 $[x_2^2+(x_2-1)^2]/2$，解为 $x_2=1/2$，目标为 $1/4$。真正的最优解是 $(0,1/2)$。第一坐标碰到边界后，第二坐标也必须重新调整，不能保持原无约束值不动。
\end{solve}

\section{正则化与偏好编码}

正则化（regularization）在数据损失之外加入对解的偏好。给定 $\alpha>0$，两个典型模型为
\begin{equation}
 \min_x\frac12\|Ax-b\|_2^2+\frac\alpha2\|x\|_2^2,
 \qquad
 \min_x\frac12\|Ax-b\|_2^2+\alpha\|x\|_1.
 \label{eq:co-regularized-fitting}
\end{equation}
前者为岭回归（ridge regression），其唯一解为 $(A^{\mathsf T}A+\alpha I)^{-1}A^{\mathsf T}b$。后者为 Lasso，最优性条件为
\[
 0\in A^{\mathsf T}(Ax-b)+\alpha\partial\|x\|_1.
\]
因此当最优分量 $x_j=0$ 时，相应残差相关量可在 $[-\alpha,\alpha]$ 内变化；这解释了绝对值正则项产生零分量的机制，但不保证任意数据都能恢复真实稀疏结构。

岭回归的稳定作用可以用奇异值分解看清。若 $A=U\Sigma V^{\mathsf T}$ 的正奇异值为 $\sigma_1,\ldots,\sigma_r$，相应左右奇异向量为 $u_j,v_j$，则
\[
 x_\alpha=\sum_{j=1}^r
 \frac{\sigma_j}{\sigma_j^2+\alpha}(u_j^{\mathsf T}b)v_j.
\]
未正则化的最小范数解使用系数 $1/\sigma_j$，很小的奇异值可能显著放大数据扰动；岭回归把它换成 $\sigma_j/(\sigma_j^2+\alpha)$，从而压低这些方向的响应。代价是即使数据没有噪声，也通常不再精确拟合。正则化改善稳定性与引入偏差，是同一次模型选择的两个方面。

\begin{example}[软阈值]
当 $A=I$ 时，Lasso 可逐坐标求解 $\min_x(x-b)^2/2+\alpha|x|$。在正半轴驻点为 $b-\alpha$，负半轴驻点为 $b+\alpha$；零点最优当且仅当 $|b|\leq\alpha$。故
\[
 x^\star=\operatorname{sign}(b)\max\{|b|-\alpha,0\},
\]
其中 $\operatorname{sign}(0)=0$，称为软阈值（soft thresholding）。对 $b=(3,1/2)$、$\alpha=1$，解为 $(2,0)$。
\end{example}

若 $x_\alpha$ 最小化 $L(x)+\alpha R(x)$，则它也最小化约束问题 $\min L(x)$、$R(x)\leq R(x_\alpha)$，否则可找到更小的惩罚目标。反向对应依赖乘子存在性，且参数不必一一对应。由两个权重下的最优性不等式相加可知，随 $\alpha$ 增大，选取的最优解的正则项值不增，数据损失值不减。

具体地，设 $0<\alpha<\beta$ 且两问题均有有限目标的最优解，则
\[
 L(x_\alpha)+\alpha R(x_\alpha)\leq L(x_\beta)+\alpha R(x_\beta),
 \qquad
 L(x_\beta)+\beta R(x_\beta)\leq L(x_\alpha)+\beta R(x_\alpha).
\]
相加得到 $(\beta-\alpha)[R(x_\beta)-R(x_\alpha)]\leq0$，再代回第一式得 $L(x_\alpha)\leq L(x_\beta)$。这个结论只描述训练模型内部的权衡，并不说明未见数据上的误差一定怎样变化。实际选取 $\alpha$ 时，可用独立验证数据比较预测损失；用来选择参数的数据与最终评价数据应在概念上区分。

\section{鲁棒损失与不确定性集合}

Huber 损失在小误差处为平方、大误差处为线性。给定 $\delta>0$，定义
\[
 h_\delta(t)=
 \begin{cases}t^2/2,&|t|\leq\delta,\\
 \delta|t|-\delta^2/2,&|t|>\delta.
 \end{cases}
\]
其导数为截断到 $[-\delta,\delta]$ 的 $t$，连续且单调不减，故损失凸。还可写成 $h_\delta(t)=\min_v\{(t-v)^2/2+\delta|v|\}$，由软阈值代入验证。这使大残差的边际惩罚不再无限增长。

辅助变量 $v$ 可以解释为由异常因素承担的那部分残差，小残差留在平方项里，大残差则部分转交给绝对值项。于是 Huber 拟合等价于联合凸问题
\[
 \min_{x,v}\ \frac12\|Ax-b-v\|_2^2+\delta\|v\|_1.
\]
当残差绝对值不超过 $\delta$ 时，最优 $v_i=0$；超过时，$v_i$ 承担超出的部分。这个表示既推导了损失的分段公式，也给出了另一种可实现的求解方式。

鲁棒优化（robust optimization）则要求约束对指定不确定性集合内的所有参数成立。设约束系数为 $a=\bar a+Pu$，$\|u\|_2\leq\rho$，其中 $P$ 为已知矩阵、$\rho\geq0$。条件 $a^{\mathsf T}x\leq b$ 对所有这样的 $u$ 成立，当且仅当
\[
 \bar a^{\mathsf T}x+\rho\|P^{\mathsf T}x\|_2\leq b.
\]
因为 $\sup_{\|u\|_2\leq\rho}u^{\mathsf T}P^{\mathsf T}x=\rho\|P^{\mathsf T}x\|_2$，上界由柯西--施瓦茨不等式给出，平行方向取到。所得为 SOCP 约束；若不确定集改为 $\|u\|_\infty\leq\rho$，则对应项为 $\rho\|P^{\mathsf T}x\|_1$。

最坏情形目标 $\sup_{u\in U}F(x,u)$ 在每个 $F(\cdot,u)$ 凸时仍凸。保证的范围由集合 $U$ 决定；抗异常损失、平均风险和最坏情形约束表达的是不同假设。

\section{函数拟合、平滑与插值}

给定固定基函数 $\phi_1,\ldots,\phi_n$，令拟合函数为 $g(t)=\sum_jx_j\phi_j(t)$。在观测位置 $t_i$ 上有 $g(t_i)=(Ax)_i$，其中 $A_{ij}=\phi_j(t_i)$，所以损失关于系数 $x$ 的凸性与基函数关于 $t$ 是否凸没有直接关系。

对网格上的函数值 $x_1,\ldots,x_n$，一阶差分矩阵 $D$ 满足 $(Dx)_i=x_{i+1}-x_i$。平方平滑项 $\|Dx\|_2^2$ 惩罚快速变化；总变差（total variation）$\|Dx\|_1$ 倾向产生较少的跳变。固定网格 $t_1<\cdots<t_n$ 上，分段线性插值函数凸当且仅当相邻斜率满足
\[
 \frac{x_{i+1}-x_i}{t_{i+1}-t_i}
 \leq\frac{x_{i+2}-x_{i+1}}{t_{i+2}-t_{i+1}},\quad i=1,\ldots,n-2.
\]
这些是关于函数值的线性不等式。对于任意其他函数族，仅在采样点检验形状通常不能认证整个连续区间上的性质。

精确插值要求 $Ax=b$，可能在噪声下放大不稳定性；逼近允许残差并通过损失和正则项进行权衡。两者首先是不同的模型选择。

\begin{example}[由单调性约束修正一组观测]
在三个依次增大的位置观测到 $b=(3,1,2)$，认为真实数值应当不减，求 $\min_x\|x-b\|_2^2/2$，满足 $x_1\leq x_2\leq x_3$。
\end{example}
\begin{solve}
前两个观测的顺序冲突，把它们调整为共同均值 $2$，得到候选点 $x=(2,2,2)$。对约束 $x_1-x_2\leq0$、$x_2-x_3\leq0$ 配乘子 $\lambda_1,\lambda_2$，驻点式为
\[
 x-b+(\lambda_1,-\lambda_1+\lambda_2,-\lambda_2)^{\mathsf T}=0.
\]
取 $(\lambda_1,\lambda_2)=(1,0)$ 即满足全部 KKT 条件，因此该候选点最优；平方目标严格凸，解唯一。第二条约束虽然活跃，乘子却为零，这也具体展示了活跃性与正乘子的区别。
\end{solve}

\section{综合例题与学习检查}
\begin{enumerate}
 \item 对观测 $0,0,M$，分别求三种范数下的最佳常数，分析 $M\to+\infty$ 时的变化。答案为 $M/3,0,M/2$。
 \item 推导带非负约束的岭回归 KKT 条件。说明为何 $A=I$ 时可以逐坐标求解，而一般 $A$ 不行。
 \item 给定线性拟合的椭球系数不确定性，写出其 SOCP 鲁棒对应形式，并核对矩阵维度。
 \item 在等距网格上拟合一列数据，写出“平方误差加总变差”和“平方误差加凸性约束”两个模型。前者偏好分段常数，后者限制斜率单调，两种模型不可互换。
 \item 推导练习：对 $A=I$、$b=(3,1/2)$，写出岭回归与 Lasso 的解随 $\alpha>0$ 变化的公式，比较哪些坐标会严格变成零。答案分别为 $b/(1+\alpha)$ 与逐坐标软阈值。
\end{enumerate}

\chapter{统计估计与实验设计}
\label{ch:co-statistics}

统计模型决定数据的分布及未知参数，优化决定按某个准则选择哪一个参数值。一个估计问题是否凸，取决于参数化、参数域和目标函数；统计一致性、偏差与泛化性质则需要另外的概率分析。

例如投掷一枚硬币十次看到七次正面，统计建模首先决定是否把十次结果看成具有同一成功概率的独立伯努利观测；只有接受这个模型，才有下面的似然函数。优化随后回答“哪个参数让已观察到的数据具有最大似然”。计算出 $0.7$ 并不表示硬币的真实正面概率被精确确定，抽样误差仍需置信区间或后验分布等工具描述。

\section{参数估计与最大似然}

设随机观测 $Z_1,\ldots,Z_N$ 在参数 $\theta\in\Theta\subseteq\mathbb R^d$ 下独立，其观测值为 $z_1,\ldots,z_N$，概率质量函数或密度为 $p_\theta$。最大似然估计（maximum likelihood estimation，MLE）最大化观测数据的似然，等价于
\begin{equation}
 \widehat\theta\in\mathop{\mathrm{arg\,min}}_{\theta\in\Theta}
 \left\{-\sum_{i=1}^N\log p_\theta(z_i)\right\}.
 \label{eq:co-mle}
\end{equation}
零概率对应负对数似然 $+\infty$。当 $\Theta$ 凸且各负对数似然关于 $\theta$ 凸时，这是凸问题。若观测不独立，应使用联合密度，不能直接写成上述乘积或求和。

对数变换有两个用途：它严格递增，因此不改变最大化的位置；它把独立观测的概率乘积变为求和，便于求导。似然是把数据固定后视为参数的函数，不是参数本身的概率分布。连续观测使用的是密度值，密度可以大于一，也不能把单个点的密度当作该点的概率。

\begin{example}[伯努利参数]
独立观测 $z_i\in\{0,1\}$，成功概率为 $q\in[0,1]$，记 $S=\sum_i z_i$。负对数似然为 $-S\log q-(N-S)\log(1-q)$。当 $0<S<N$ 时，令导数 $-S/q+(N-S)/(1-q)=0$，得到唯一解 $\widehat q=S/N$。当 $S=0$ 或 $S=N$ 时，解分别为端点 $0,1$，相应零计数项按连续延拓处理。
\end{example}

指数族在自然参数 $\theta$ 下写成 $p_\theta(z)=h(z)\exp(\theta^{\mathsf T}T(z)-A(\theta))$。对数配分函数 $A$ 凸，所以负对数似然是凸函数加仿射项。常用例子是逻辑回归：给定特征 $a_i$ 和标签 $y_i\in\{0,1\}$，目标为
\[
 \sum_i\bigl[\log(1+e^{a_i^{\mathsf T}\theta})-y_i a_i^{\mathsf T}\theta\bigr].
\]
海塞矩阵为 $\sum_iq_i(1-q_i)a_i a_i^{\mathsf T}\succeq0$，其中 $q_i=(1+e^{-a_i^{\mathsf T}\theta})^{-1}$。数据完全可分时，无正则项的最优值可能只在参数范数趋于无穷时逼近，凸性并不保证有限最优参数。

上式可从伯努利模型逐步得到。令 $s_i=a_i^{\mathsf T}\theta$，通过 $q_i=1/(1+e^{-s_i})$ 把任意实数分数转为概率，则单点似然为 $q_i^{y_i}(1-q_i)^{1-y_i}$。取负对数并化简，得到 $\log(1+e^{s_i})-y_i s_i$；对 $\theta$ 求导得到 $(q_i-y_i)a_i$。梯度因此由“预测概率减观测标签”乘以特征向量组成。

\begin{example}[只有截距的逻辑回归]
十个二元标签中有七个一、三个零，所有特征都为 $a_i=1$，求逻辑回归参数。
\end{example}
\begin{solve}
目标为 $F(\theta)=10\log(1+e^\theta)-7\theta$，梯度为 $10q-7$，所以 $q^\star=0.7$、$\theta^\star=\log(7/3)$。概率参数 $q$ 与对数优势参数 $\theta=\log(q/(1-q))$ 表示同一个模型，却使用不同坐标。若十个标签全为一，目标变成 $10\log(1+e^{-\theta})$，下确界为零，只能在 $\theta\to+\infty$ 时逼近；改回闭区间上的概率参数则有端点解 $q=1$。域的选择决定最优参数是否存在。
\end{solve}

指数族的配分函数也有类似的凸性原因。对有限支持，若 $Z(\theta)=\sum_z h(z)e^{\theta^{\mathsf T}T(z)}$、$A(\theta)=\log Z(\theta)$，直接求导得 $\nabla A(\theta)=\mathbb E_\theta T(Z)$，$\nabla^2 A(\theta)=\operatorname{Cov}_\theta(T(Z))\succeq0$。协方差半正定解释了自然参数下的凸性；连续支持时须先保证积分存在，并验证允许把微分移入积分号的条件。

\section{最大后验估计与正则化}

给定先验密度 $\pi(\theta)$，最大后验估计（maximum a posteriori estimation，MAP）最小化
\[
 -\sum_i\log p_\theta(z_i)-\log\pi(\theta),
\]
其中省略与 $\theta$ 无关的归一化常数。若似然和先验关于参数均对数凹，且共同支持集凸，则负对数后验凸。高斯先验产生平方正则项，独立拉普拉斯先验产生绝对值正则项；尺度决定对应系数。

\begin{example}[高斯均值的 MAP]
设 $Z_i\mid\theta\sim N(\theta,\sigma^2)$ 独立，$\sigma^2>0$ 已知；先验为 $\theta\sim N(m_0,\tau^2)$，$\tau^2>0$。负对数后验除常数外为
\[
 \frac1{2\sigma^2}\sum_i(z_i-\theta)^2+\frac1{2\tau^2}(\theta-m_0)^2.
\]
其唯一最小点为
\[
 \widehat\theta_{\mathrm{MAP}}
 =\frac{(N/\sigma^2)\bar z+(1/\tau^2)m_0}{N/\sigma^2+1/\tau^2},
 \qquad \bar z=\frac1N\sum_i z_i.
\]
这按精度对样本均值和先验均值加权。本例后验为高斯，众数与均值相同；一般 MAP 只给出后验众数，不能替代整个后验分布。
\end{example}

\section{有限支持上的分布估计}

在给定有限支持 $\{z_1,\ldots,z_s\}$ 上，以概率向量 $q\in\Delta_s$ 为变量。已知矩条件可写为 $Bq=c$，例如均值条件 $\sum_jq_jz_j=\mu$。熵（entropy）为 $H(q)=-\sum_jq_j\log q_j$，约定 $0\log0=0$。由于 $t\log t$ 在 $t>0$ 上凸且可连续延拓到零，最大熵模型
\[
 \max_{q\geq0}\ H(q),\qquad \mathbf1^{\mathsf T}q=1,\quad Bq=c
\]
是凸优化。若可行集非空，它紧且熵连续，因此最优解存在；熵严格凹，故解唯一。

无矩约束时，对称性和 Jensen 不等式给出均匀分布 $q_j=1/s$。若最优解全部为正，对负熵写 KKT 可得 $q_j$ 与 $\exp(-b_j^{\mathsf T}\nu)$ 成比例，其中 $b_j$ 为 $B$ 的第 $j$ 列，乘子 $\nu$ 由矩约束确定。

给定参考分布 $r$，相对熵（relative entropy）为 $D(q\|r)=\sum_jq_j\log(q_j/r_j)$。当 $q_j>0,r_j=0$ 时定义该项为 $+\infty$；$q_j=0$ 时定义为零。固定 $r$ 时它关于 $q$ 凸。有限支持上的优化是有限维问题，将连续分布离散化后得到的解只对该离散模型直接最优。

相对熵虽常用于衡量分布差别，却不是对称距离。例如一般有 $D(q\|r)\ne D(r\|q)$。其非负性可由 $\log u\leq u-1$ 推出：对所有 $q_j>0$ 的项取 $u=r_j/q_j$，得到
\[
 D(q\|r)\geq\sum_{j:q_j>0}(q_j-r_j)\geq0.
\]
若某个 $q_j>0$ 而 $r_j=0$，则非负性由无穷值约定直接成立。有限值时等号最终要求 $q=r$。这既说明参考分布本身代价最低，也说明支持集的零概率不能随意忽略。

\section{检测器设计与假设检验}

考虑两个假设 $H_0,H_1$，观测只能取有限值 $z_1,\ldots,z_s$，相应已知概率为 $p_{0\ell},p_{1\ell}$。令 $t_\ell\in[0,1]$ 表示观测 $z_\ell$ 后判为 $H_1$ 的概率。$t_\ell\in\{0,1\}$ 对应确定性检测器，允许区间值则为随机检测器（randomized detector）。

虚警率为 $\alpha(t)=\sum_\ell p_{0\ell}t_\ell$，漏检率为 $\beta(t)=\sum_\ell p_{1\ell}(1-t_\ell)$。给定虚警上限 $a\in[0,1]$，最小化 $\beta(t)$，满足 $\alpha(t)\leq a$、$0\leq t\leq1$，是 LP。

对虚警约束引入乘子 $\lambda\geq0$，拉格朗日目标中 $t_\ell$ 的系数为 $\lambda p_{0\ell}-p_{1\ell}$。故当该系数为负时选 $t_\ell=1$，为正时选零，相等时可随机化。对于 $p_{0\ell}>0$，这就是按似然比 $p_{1\ell}/p_{0\ell}$ 与阈值比较；阈值由虚警限制决定。优化变量是决策规则，输入分布是模型给定的数据。

\begin{example}[随机化为什么会改善检测器]
观测只有两个取值，两假设下的分布分别为 $p_0=(0.8,0.2)$、$p_1=(0.2,0.8)$，允许虚警率不超过 $0.1$。
\end{example}
\begin{solve}
确定性检测器若在任一观测上判为 $H_1$，虚警至少为 $0.2$，因此唯一可行选择是永远判为 $H_0$，漏检率为一。允许随机化后取 $t=(0,1/2)$，虚警为 $0.1$、漏检为 $0.6$。这是最优的，因为检测成功率满足
\[
 0.2t_1+0.8t_2\leq4(0.8t_1+0.2t_2)\leq0.4,
\]
而该规则恰好达到 $0.4$。随机化将离散规则的可行集合扩展为凸集，并在这个例子中带来严格改善。
\end{solve}

\section{概率界与优化表达}

对非负随机变量 $Y$ 和 $a>0$，$\mathbb EY\geq a\mathbb P(Y\geq a)$ 给出 Markov 不等式。若 $Z$ 的均值为 $\mu$、方差为 $\sigma^2<\infty$，令 $Y=(Z-\mu)^2$，得切比雪夫界（Chebyshev bound）
\[
 \mathbb P(|Z-\mu|\geq a)\leq\sigma^2/a^2.
\]
若 $M_Z(t)=\mathbb E e^{tZ}$ 在某个 $t>0$ 处有限，再对 $e^{tZ}$ 使用 Markov 不等式，得到 Chernoff 界
\[
 \mathbb P(Z\geq a)\leq\inf_{t\geq0:\,M_Z(t)<\infty}
       \exp\bigl(\log M_Z(t)-ta\bigr).
\]
对数矩母函数在有限性区间上凸；可由 Hölder 不等式得到。因此选取最紧的这一类上界是一个单变量凸问题。所得仍是概率的上界，不一定等于真实概率。

若只知道有限支持上的若干矩，最大化指定事件的概率等于最大化 $\sum_{j:z_j\in E}q_j$，满足概率和矩约束，这是 LP。对连续支持，变量变为概率分布本身；有限维近似或矩方法给出的界应与原分布问题区分。

\section{实验设计与信息矩阵}

在线性观测模型 $Y=a_i^{\mathsf T}\theta+\varepsilon$ 中，假设噪声独立、均值为零且方差为已知 $\sigma^2$；若解释为 Fisher 信息，再假设噪声为高斯。候选测量向量为 $a_1,\ldots,a_s\in\mathbb R^d$，连续设计权重满足 $w\in\Delta_s$，定义归一化信息矩阵
\[
 M(w)=\sum_iw_i a_i a_i^{\mathsf T}.
\]
总测量数为 $N$、且各方向次数为 $Nw_i$ 时，总信息为 $NM(w)/\sigma^2$。在 $M(w)\succ0$ 的域上，D-最优设计最小化 $-\log\det M(w)$；A-最优设计最小化 $\operatorname{tr}(M(w)^{-1})$；E-最优设计最大化最小特征值，可写为最大化 $t$，满足 $M(w)\succeq tI$。

这些准则来自估计误差的方向性。将实际测量向量组成设计矩阵 $A$，若满列秩且噪声协方差为 $\sigma^2I$，最小二乘估计满足
\[
 \widehat\theta-\theta=(A^{\mathsf T}A)^{-1}A^{\mathsf T}\varepsilon,
 \qquad
 \operatorname{Cov}(\widehat\theta)=\sigma^2(A^{\mathsf T}A)^{-1}.
\]
当 $A^{\mathsf T}A=NM(w)$ 时，信息矩阵某方向上的特征值越小，该方向上的估计方差越大。因此只测量一个方向，即使次数很多，也可能完全无法识别与它正交的参数。设计问题调整的是未来采集数据的比例，区别于拿到数据之后再拟合参数。

A-最优设计的凸性可由 Schur 补看出：引入对称矩阵 $Z$，最小化 $\operatorname{tr}Z$，满足
\[
 \begin{bmatrix}M(w)&I\\I&Z\end{bmatrix}\succeq0,
\]
即要求 $Z\succeq M(w)^{-1}$。这些设计分别控制误差椭球体积、平均方差和最不利方向信息；连续权重是整数次数的松弛，取整后须重新检查矩阵的秩与目标值。

\section{例题与学习检查}
\begin{example}[两个正交测量方向]
设 $a_1=(1,0)^{\mathsf T}$、$a_2=(0,1)^{\mathsf T}$，则 $M(w)=\operatorname{diag}(w,1-w)$，$0<w<1$。D-最优设计最大化 $w(1-w)$，A-最优设计最小化 $1/w+1/(1-w)$，E-最优设计最大化 $\min\{w,1-w\}$；三者均在 $w=1/2$ 取最优。这一重合来自该例的对称性，并非一般结论。
\end{example}
\begin{enumerate}
 \item 推导独立 Poisson 观测的均值参数最大似然估计，分别在均值参数和对数均值参数下判断凸性。
 \item 在三点支持 $\{0,1,2\}$ 上最大化熵，要求均值为 $1$。验证均匀分布最优，并说明只给定归一化与均值时为何解仍唯一。
 \item 对 $Z\sim N(0,1)$，用 $\log M_Z(t)=t^2/2$ 优化 Chernoff 界，求 $a\geq0$ 时的结果 $e^{-a^2/2}$。
 \item 综合练习：给只有截距的逻辑回归加入 $\alpha\theta^2/2$，$\alpha>0$。证明无论二元标签怎样分布，最优参数都存在且唯一。提示：平方项给出强凸性和强制性，不需要假设样本中两类都出现。
\end{enumerate}

\chapter{几何优化与分类}
\label{ch:co-geometry}

距离、包含、分离和位置都可以转化为优化约束。建模的关键是选择合适的变量：同一个几何对象采用不同参数化，可能显露或遮蔽凸结构。

\section{向凸集投影与最优性}

对非空闭凸集 $C$，欧氏投影（Euclidean projection）定义为
\begin{equation}
 \Pi_C(z)=\mathop{\mathrm{arg\,min}}_{x\in C}\frac12\|x-z\|_2^2.
 \label{eq:co-projection}
\end{equation}
根据定理\ref{thm:co-projection-existence}，最优解唯一，故这里将最优解集合与其中唯一元素认同。由一阶最优性条件，$u=\Pi_C(z)$ 当且仅当
\[
 (z-u)^{\mathsf T}(x-u)\leq0\quad\text{对所有 }x\in C.
\]
这表示从投影点指向给定点 $z$ 的方向，与所有可行位移都成非锐角。

\begin{proposition}[投影的非扩张性]
$\|\Pi_C(z)-\Pi_C(w)\|_2\leq\|z-w\|_2$。
\end{proposition}
\begin{proof}
令 $u=\Pi_C(z)$、$v=\Pi_C(w)$，在两条投影不等式中分别取 $x=v,u$ 并相加，得 $\|u-v\|_2^2\leq(u-v)^{\mathsf T}(z-w)$。再用柯西--施瓦茨不等式；若 $u=v$ 则结论直接成立。
\end{proof}

到半空间的投影已在第\ref{ch:co-sets}章求出。点 $z$ 到球 $\{x:\|x-c\|_2\leq r\}$ 的投影在 $z$ 位于球内时为 $z$，在球外时为 $c+r(z-c)/\|z-c\|_2$。到满行秩仿射集合 $Ax=b$ 的投影为 $z-A^{\mathsf T}(AA^{\mathsf T})^{-1}(Az-b)$，可由最小范数修正求得。

\begin{example}[投影到概率单纯形]
给定 $z\in\mathbb R^n$，求使 $\|x-z\|_2^2/2$ 最小的概率向量 $x\geq0$、$\sum_i x_i=1$。
\end{example}
\begin{solve}
给等式配乘子 $\tau$，给 $-x_i\leq0$ 配乘子 $\lambda_i$。驻点条件为 $x_i-z_i+\tau-\lambda_i=0$。若 $x_i>0$，则 $\lambda_i=0$，故 $x_i=z_i-\tau$；若 $x_i=0$，则 $\lambda_i=\tau-z_i\geq0$。两种情况合并为
\begin{equation}
 x_i^\star=\max\{z_i-\tau,0\},\qquad
 \sum_i\max\{z_i-\tau,0\}=1.
 \label{eq:co-simplex-projection}
\end{equation}
最后一式的左侧连续，在正值区域严格递减，且从 $+\infty$ 下降到零，所以恰有一个解 $\tau$。可用二分或排序寻找它。例如 $z=(0.8,0.5,-0.1)$ 时，取 $\tau=0.15$ 得 $x^\star=(0.65,0.35,0)$。投影同时满足非负与总和限制；只把负数截成零后再按比例归一化，一般不是欧氏距离下的最近点。
\end{solve}

投影也说明了怎样把一个不符合约束的试探方案修正为可行方案。非扩张性保证输入的小变化不会在修正后被放大，这正是第\ref{sec:co-projected-gradient}节投影梯度法的基础。它与第\ref{ch:co-sets}章中消去坐标的线性投影不同：到一般凸集的最近点映射通常是非线性的。

\section{集合距离与分离问题}

非空凸集 $C,D$ 的距离为 $\inf_{x\in C,y\in D}\|x-y\|_2$，是关于 $(x,y)$ 的凸优化。若一个集合紧、另一个闭，取极小化序列并利用有界性可得到收敛子列，因此距离取到；不相交时该距离严格为正。

但两个闭凸集不相交时距离仍可能为零。例如 $C=\{(x,y):y\geq e^x\}$、$D=\{(x,y):y\leq0\}$，取 $(x,e^x)$ 与 $(x,0)$，令 $x\to-\infty$ 即可使距离趋零，且永不相交。

支撑函数 $\sigma_C(v)=\sup_{x\in C}v^{\mathsf T}x$ 将点到集合的距离写为
\[
 \operatorname{dist}(z,C)=\max_{\|v\|_2\leq1}\{v^{\mathsf T}z-\sigma_C(v)\}
\]
，其中 $C$ 非空闭凸。任意可行 $v$ 和 $x\in C$ 满足 $v^{\mathsf T}z-\sigma_C(v)\leq\|z-x\|_2$；若 $z\notin C$，取 $v=(z-\Pi_C(z))/\|z-\Pi_C(z)\|_2$，投影不等式保证下界取等号。若 $z\in C$，取 $v=0$。这给出了距离问题的一个对偶解释。

\section{距离矩阵、格拉姆矩阵与半正定表示}

设未知点 $p_1,\ldots,p_N\in\mathbb R^d$ 作为矩阵 $P$ 的行，格拉姆矩阵（Gram matrix）为 $G=PP^{\mathsf T}\succeq0$。平方距离满足
\[
 \|p_i-p_j\|_2^2=G_{ii}+G_{jj}-2G_{ij}.
\]
因此给定部分平方距离会形成关于 $G$ 的仿射约束。反之，若 $G\succeq0$ 且秩为 $r$，由谱分解 $G=U_r\Lambda_rU_r^{\mathsf T}$，矩阵 $P=U_r\Lambda_r^{1/2}$ 的行给出 $\mathbb R^r$ 中的一组坐标。

若要求固定低维空间，例如 $d=2$，还须加上 $\operatorname{rank}G\leq d$，这是非凸限制。删除秩约束得到 SDP 松弛，得到的点集可能需要更高维空间才能准确实现。截取最大的 $d$ 个特征值可以给出近似坐标，但应重新计算距离误差。

距离不随整体平移而改变，格拉姆矩阵却会改变，所以仅从距离通常不能确定唯一的 $G$。可加入中心化条件 $G\mathbf1=0$，对应点集质心位于原点。即使如此，坐标仍可整体旋转或反射；这些变换不改变任何距离，不应被误认为求解有误。优化中的唯一性应针对所选变量和实际可识别的对象分别讨论。

\section{极值体积椭球与不同中心}

用 $E=\{x:\|Bx+d\|_2\leq1\}$ 表示椭球，其中 $B\succ0$。其体积为单位球体积除以 $\det B$，故包围点 $a_i$ 的最小体积椭球可求解
\[
 \min_{B\succ0,d}-\log\det B,
 \qquad \|Ba_i+d\|_2\leq1\quad\text{对所有 }i.
\]
约束关于 $(B,d)$ 凸。若点集仿射张成全空间，则可得到非退化最小包围椭球；若点集低维，满维体积的下确界可能为零而不取到。

对于多面体 $C=\{x:a_i^{\mathsf T}x\leq b_i\}$，内接椭球用 $E=\{c+Bu:\|u\|_2\leq1\}$、$B\succ0$ 表示。其包含条件等价于 $a_i^{\mathsf T}c+\|Ba_i\|_2\leq b_i$，体积正比于 $\det B$，故最大化 $\log\det B$ 是凸优化意义下的凹最大化。存在满维可行椭球需多面体有内点；有界性使最大体积有限。

切比雪夫中心（Chebyshev center）选择最大内接球中心，可通过 LP 求得：最大化 $r\geq0$，满足 $a_i^{\mathsf T}c+r\|a_i\|_2\leq b_i$。解析中心（analytic center）则最小化 $-\sum_i\log(b_i-a_i^{\mathsf T}c)$，定义于严格可行域。对于有界、具有严格可行点的多面体，目标逼近边界时趋于无穷，故中心存在；不同的冗余约束表示可能改变解析中心。

椭球模型的变量选择尤其值得注意。直接把中心 $c$ 和形状矩阵 $B$ 同时放进 $\|B(x-c)\|_2\leq1$，会出现变量乘积 $Bc$；包围椭球使用独立变量 $d=-Bc$ 后，样本约束 $\|Ba_i+d\|_2\leq1$ 才显出仿射复合结构。求解后再恢复 $c=-B^{-1}d$。等价重参数化的价值就在于把被变量乘积遮住的凸性显露出来。

\section{线性分类与最大间隔}

给定样本 $a_i\in\mathbb R^n$、标签 $y_i\in\{-1,1\}$，假设两类样本都存在，分类超平面为 $w^{\mathsf T}x+\beta=0$，其中 $w\ne0$。若数据线性可分，可缩放 $w,\beta$，使 $y_i(w^{\mathsf T}a_i+\beta)\geq1$。点到超平面的距离为 $|w^{\mathsf T}a_i+\beta|/\|w\|_2$，因此最大化规范化后的间隔等价于硬间隔支持向量机（support vector machine，SVM）：
\[
 \min_{w,\beta}\frac12\|w\|_2^2,
 \qquad y_i(w^{\mathsf T}a_i+\beta)\geq1.
\]
不可分时引入松弛量 $\xi_i\geq0$，将右侧改为 $1-\xi_i$，目标加入 $C\sum_i\xi_i$，其中 $C>0$。消去松弛量就得到铰链损失 $\max\{0,1-y_i(w^{\mathsf T}a_i+\beta)\}$。

硬间隔模型的驻点条件给出 $w=\sum_i\alpha_i y_i a_i$、$\sum_i\alpha_i y_i=0$、$\alpha_i\geq0$；互补松弛为 $\alpha_i[y_i(w^{\mathsf T}a_i+\beta)-1]=0$。因此只有间隔约束活跃的样本可能有非零乘子，它们决定分类法向量，称为支持向量。

这一归一化不是对原始数据的额外假设。对有限个严格分离的样本，最小的带符号分数 $m=\min_i y_i(w^{\mathsf T}a_i+\beta)$ 为正；同时除以 $m$ 就使最小分数恰为一，几何超平面没有改变。规范化之后，两条支持平面 $w^{\mathsf T}x+\beta=\pm1$ 的距离为 $2/\|w\|_2$，所以最小化 $\|w\|_2^2/2$ 正是在最大化这个间隔。

代入驻点条件还可得到硬间隔对偶：
\[
 \max_{\alpha\geq0}\
 \sum_i\alpha_i-\frac12\left\|\sum_i\alpha_i y_i a_i\right\|_2^2,
 \qquad \sum_i\alpha_i y_i=0.
\]
对偶把优化变量从特征维度转移到样本维度，并只通过内积 $a_i^{\mathsf T}a_j$ 使用特征。软间隔模型中，对 $\xi_i$ 的驻点条件进一步产生 $0\leq\alpha_i\leq C$；上界反映允许违反间隔时所支付的惩罚。这里仅解释线性模型，统计泛化能力仍取决于数据与模型假设。

\section{设施定位、网络布局与平面布置}

给定需求点 $a_i$ 和非负权重 $w_i$，设施位置 $x$ 可由 $\min_x\sum_iw_i\|x-a_i\|_2$ 确定。引入 $t_i\geq\|x-a_i\|_2$ 即为 SOCP；区域约束若为凸集，模型仍凸。多个设施且连接关系固定时，代价 $\sum_{(i,j)}w_{ij}\|x_i-x_j\|_2$ 也凸。

若同时决定哪个设施服务哪个客户，会出现离散指派；这部分通常不能直接保留在纯凸模型中。矩形非重叠也包含“左、右、上、下”中的离散选择。只有预先固定相对关系，或采用其他具有明确保证的松弛，才可将相应位置限制写成单组凸约束。

\section{例题与学习检查}
\begin{example}[一维最大间隔]
负类点为 $-2$，正类点为 $2$。硬间隔约束为 $2w-\beta\geq1$、$2w+\beta\geq1$，相加得 $w\geq1/2$。故最小化 $w^2/2$ 时取 $w=1/2$、$\beta=0$，分类边界为原点，两侧几何间隔均为 $2$。
\end{example}
\begin{enumerate}
 \item 求 $(3,4)$ 到单位闭球的投影，验证距离为 $4$，投影为 $(3/5,4/5)$。
 \item 推导盒状区域 $l_i\leq x_i\leq u_i$ 上的欧氏投影。答案为逐坐标截断。
 \item 比较区间 $[0,3]$ 的最大内接球中心与解析中心；再添加冗余约束 $x<4$，检查解析中心是否改变。
 \item 给定三点的部分距离，写出 $3\times3$ 格拉姆矩阵的 SDP 可行性模型，并指出若要求点都在直线上须补充什么条件。
 \item 算法准备：按式\eqref{eq:co-simplex-projection}求 $z=(2,1,-1)$ 到三维概率单纯形的投影，写出全部 KKT 乘子。核对结果为 $(1,0,0)$，阈值为 $\tau=1$。
\end{enumerate}

\part{凸优化算法}

\chapter{梯度方法与牛顿方法}
\label{ch:co-unconstrained}

本章从无约束问题 $\min_x f(x)$ 出发，再介绍处理简单约束和非光滑项的一阶方法，最后讨论牛顿法。除专门讨论非光滑问题的节外，函数凸且可微；出现 $x^\star$ 时假设最优解存在，以 $f^\star=f(x^\star)$ 记最优值。算法产生序列 $x_0,x_1,\ldots$，必须区分目标误差、点误差和计算成本。

解析求解与数值求解回答同一个最优化问题，只是实现方式不同。例如最小二乘可写出正规方程，但高维时仍要数值解线性系统；其他模型甚至没有简洁的闭式方程解。迭代算法通过反复解决容易的局部问题，逐步改善当前方案。下面每种算法都应从“本步实际最小化了哪个简单模型”来理解。

\section{算法问题、误差与基本假设}

目标误差为 $f(x_k)-f^\star$，点误差为 $\|x_k-x^\star\|_2$，梯度范数 $\|\nabla f(x_k)\|_2$ 则直接可计算。一般情况下，函数值接近最优不保证点接近某个指定最优解；若函数为 $\mu$-强凸，则
\[
 \frac\mu2\|x-x^\star\|_2^2\leq f(x)-f^\star
 \leq\frac1{2\mu}\|\nabla f(x)\|_2^2.
\]
左侧来自最优点处的强凸下界；右侧对 $f(y)\geq f(x)+\nabla f(x)^{\mathsf T}(y-x)+\mu\|y-x\|_2^2/2$ 关于 $y$ 取下确界得到。

例如 $f(x_1,x_2)=x_1^2$ 的最优解集是整条直线 $x_1=0$。点 $(0,10^6)$ 已经最优，却离指定最优点 $(0,0)$ 很远，所以函数值收敛不能自动解释为靠近某个固定解。即使解唯一，函数过于平坦也会使小梯度缺乏精确意义；强凸参数正是把梯度、目标误差与点误差联系起来所需的定量信息。

\begin{lemma}[下降引理]
若 $f$ 在包含线段 $[x,y]$ 的开凸域上可微，且梯度以常数 $L$ 利普希茨连续，则
\begin{equation}
 f(y)\leq f(x)+\nabla f(x)^{\mathsf T}(y-x)+\frac L2\|y-x\|_2^2.
 \label{eq:co-descent-lemma}
\end{equation}
\end{lemma}
\begin{proof}
令 $d=y-x$。由微积分基本定理，$f(y)-f(x)-\nabla f(x)^{\mathsf T}d=\int_0^1[\nabla f(x+td)-\nabla f(x)]^{\mathsf T}d\,dt$，其上界为 $\int_0^1Lt\|d\|_2^2\,dt=L\|d\|_2^2/2$。
\end{proof}

\section{下降方向与线搜索}

统一更新为 $x_{k+1}=x_k+t_kd_k$，其中 $d_k$ 为方向，$t_k>0$ 为步长。当 $\nabla f(x_k)^{\mathsf T}d_k<0$ 时，称 $d_k$ 为下降方向（descent direction）。可微性保证充分小的正步长使目标下降，但过大步长可能越过低谷或离开定义域。

Armijo 回溯线搜索（backtracking line search）取参数 $0<\alpha<1/2$、$0<\beta<1$，先令 $t=1$；若试探点不在定义域内，或
\[
 f(x+td)>f(x)+\alpha t\nabla f(x)^{\mathsf T}d,
\]
就令 $t\leftarrow\beta t$，直到满足条件。若当前点位于开定义域且 $d$ 为下降方向，则足够小步长必满足条件，因为一阶展开中的余项为 $o(t)$。精确线搜索直接最小化射线上的目标，通常比回溯更贵。

例如 $f(x)=x^2/2$，在 $x=1$ 处沿 $d=-1$ 试探。真实变化为 $f(1-t)-f(1)=-t+t^2/2$，所以目标下降要求 $0<t<2$；Armijo 条件进一步要求 $t\leq2(1-\alpha)$。它不是只检查目标是否减少，而是要求减少量至少占一阶预测下降量的一定比例，避免接受过于微弱的改善。

\section{梯度下降与收敛速度}

\label{sec:co-gradient}

梯度指向局部增长最快的欧氏方向，因此其相反方向自然适合减小目标。但只最小化线性模型 $f(x)+\nabla f(x)^{\mathsf T}(y-x)$ 会使非零梯度下的最小值无界。加入一个限制位移的二次项后，
\[
 \mathop{\mathrm{arg\,min}}_y
 \left\{f(x)+\nabla f(x)^{\mathsf T}(y-x)+\frac1{2t}\|y-x\|_2^2\right\}
 =x-t\nabla f(x).
\]
因此步长 $t$ 决定愿意在多大范围内相信当前的线性模型；$t$ 越小，对远距离移动的惩罚越强。

梯度下降（gradient descent）采用 $d_k=-\nabla f(x_k)$。给定 $x_0$ 和误差要求，可使用已知 $L>0$ 的固定步长 $1/L$，或每步回溯；更新后重复计算梯度，达到适用的停止条件时返回当前点。

\begin{example}[手算一次完整的梯度迭代]
对 $f(x)=(x-3)^2/2$，从 $x_0=0$ 出发，采用步长 $t=1/2$。
\end{example}
\begin{solve}
梯度为 $x-3$，递推为 $x_{k+1}=(x_k+3)/2$。前几步为 $0,3/2,9/4,21/8$，归纳得
\[
 x_k=3(1-2^{-k}),\qquad
 f(x_k)-f^\star=\frac92\,4^{-k}.
\]
每一步都把点误差减半，把目标误差变成原来的四分之一。若换成步长 $t$，则 $x_k-3=(1-t)^k(x_0-3)$；$0<t<2$ 时收敛，$t=2$ 时一般在两点间振荡，$t>2$ 时误差增大。方向正确并不能补救任意过大的步长。
\end{solve}

\begin{theorem}[光滑凸情形的目标误差界]
设 $f:\mathbb R^n\to\mathbb R$ 凸，梯度以常数 $L>0$ 利普希茨连续，且存在最优解。令 $x_{k+1}=x_k-\nabla f(x_k)/L$，则对 $k\geq1$，
\begin{equation}
 f(x_k)-f^\star\leq\frac{L\|x_0-x^\star\|_2^2}{2k}.
 \label{eq:co-gradient-rate}
\end{equation}
\end{theorem}
\begin{proof}
记 $g_k=\nabla f(x_k)$。下降引理给出 $f(x_{k+1})\leq f(x_k)-\|g_k\|_2^2/(2L)$，故目标值不增。再用凸性 $f(x_k)-f^\star\leq g_k^{\mathsf T}(x_k-x^\star)$，得
\[
 f(x_{k+1})-f^\star
 \leq\frac L2\bigl(\|x_k-x^\star\|_2^2-\|x_{k+1}-x^\star\|_2^2\bigr).
\]
从 $0$ 至 $k-1$ 求和，右侧望远镜消去；左侧至少为 $k[f(x_k)-f^\star]$，于是得到所述界。
\end{proof}

若再有 $\mu$-强凸性，则前节梯度误差界与下降引理给出
\[
 f(x_{k+1})-f^\star\leq(1-\mu/L)[f(x_k)-f^\star].
\]
因此目标误差按几何比例衰减，称为线性收敛（linear convergence）。达到给定误差所需迭代次数与条件比 $L/\mu$ 及精度对数相关；一般光滑凸情形只有式\eqref{eq:co-gradient-rate}的 $O(1/k)$ 目标误差界。这里 $O(1/k)$ 描述误差随迭代次数的衰减阶；反解误差要求后，所需迭代次数为 $O(1/\varepsilon)$，尚不包含每次梯度计算的成本。

\section{投影梯度法：处理简单凸约束}
\label{sec:co-projected-gradient}

设 $C\subseteq\mathbb R^n$ 非空闭凸，目标 $f:\mathbb R^n\to\mathbb R$ 凸且梯度以 $L>0$ 利普希茨连续。直接作梯度步可能离开 $C$，可以先试探，再把结果投影回来：
\begin{equation}
 x_{k+1}=\Pi_C\bigl(x_k-t\nabla f(x_k)\bigr),\qquad x_0\in C.
 \label{eq:co-projected-gradient}
\end{equation}
这称为投影梯度法（projected gradient method）。配方可知，它恰好在 $y\in C$ 上最小化前节的二次局部模型。盒约束、球和概率单纯形的投影较容易计算，所以该方法特别适合这些模型；若投影本身需要解一个大规模优化问题，每步就未必便宜。

\begin{proposition}[投影梯度法的下降与误差界]
在上述条件下，假设最优解 $x^\star$ 存在，取固定 $0<t\leq1/L$。则目标值不增，并对 $k\geq1$ 有
\[
 f(x_k)-f^\star\leq\frac{\|x_0-x^\star\|_2^2}{2tk}.
\]
\end{proposition}
\begin{proof}
记 $x^+=\Pi_C(x-t\nabla f(x))$。投影的一阶最优性给出
\[
 \left(\nabla f(x)+\frac{x^+-x}{t}\right)^{\mathsf T}(y-x^+)\geq0,
 \qquad y\in C.
\]
将此式与 $f(x^+)\leq f(x)+\nabla f(x)^{\mathsf T}(x^+-x)+\|x^+-x\|_2^2/(2t)$、$f(y)\geq f(x)+\nabla f(x)^{\mathsf T}(y-x)$ 相结合，再展开平方，得到
\begin{equation}
 f(x^+)-f(y)\leq\frac1{2t}
 \bigl(\|x-y\|_2^2-\|x^+-y\|_2^2\bigr).
 \label{eq:co-projected-three-point}
\end{equation}
取 $y=x$ 得目标不增；取 $y=x^\star$ 并对各步求和，再用目标单调性，即得所述界。
\end{proof}

约束最优点的梯度可以不为零，因此停止时不能仍机械地检查 $\|\nabla f(x)\|_2$。定义投影梯度映射
\[
 G_t(x)=\frac{x-\Pi_C(x-t\nabla f(x))}{t}.
\]
由投影最优性，$G_t(x)=0$ 当且仅当式\eqref{eq:co-first-order-optimality}成立。小的 $\|G_t(x)\|_2$ 表示接近这一固定点条件；将它转成特定目标精度仍需额外误差界或对偶证书。

\begin{example}[边界最优点的梯度不消失]
最小化 $f(x)=(x-2)^2/2$，满足 $0\leq x\leq1$。从 $x_0=0$ 采用 $t=1$，先得到无约束试探点 $2$，再投影为 $x_1=1$，恰好达到最优点。此处 $f'(1)=-1$，但 $G_1(1)=1-\Pi_{[0,1]}(2)=0$。约束阻止了继续向右移动，算法应停止而不是等待普通梯度变成零。
\end{example}

\section{近端梯度法：处理可分离的非光滑项}
\label{sec:co-proximal-gradient}

Lasso 的绝对值项在零点不可微，普通梯度法不能直接套用。考虑复合目标
\[
 F(x)=h(x)+R(x),
\]
其中 $h:\mathbb R^n\to\mathbb R$ 凸且梯度以 $L>0$ 利普希茨连续，$R$ 为适当闭凸函数，可以取 $+\infty$ 来编码约束。对 $t>0$，定义近端算子（proximal operator）
\[
 \operatorname{prox}_{tR}(v)
 =\mathop{\mathrm{arg\,min}}_y
 \left\{R(y)+\frac1{2t}\|y-v\|_2^2\right\}.
\]
该最优点存在且唯一：适当闭凸函数在有限维空间中有一个仿射下界，这可由上图分离得到；加上正二次项后次水平集有界，闭性给出存在性，强凸性给出唯一性。当 $R=I_C$ 时，近端算子就是投影；当 $R(x)=\alpha\|x\|_1$ 时，就是逐坐标阈值为 $t\alpha$ 的软阈值。

近端梯度法（proximal gradient method）只线性化光滑部分，保留非光滑部分的完整结构：
\begin{equation}
 x_{k+1}=\operatorname{prox}_{tR}\bigl(x_k-t\nabla h(x_k)\bigr).
 \label{eq:co-proximal-gradient}
\end{equation}
它等价于最小化
\[
 Q_x(y)=h(x)+\nabla h(x)^{\mathsf T}(y-x)
             +\frac1{2t}\|y-x\|_2^2+R(y).
\]
当 $t\leq1/L$ 时，$Q_x$ 在每一点都不小于 $F$，而在 $y=x$ 处与 $F$ 相等，因此每次精确最小化这个模型都会使原目标下降。

\begin{proposition}[近端梯度法的基本收敛界]
假设上述 $F$ 有最优解 $x^\star$，$x_0\in\operatorname{dom}R$，固定步长 $0<t\leq1/L$，并精确计算近端点。则
\[
 F(x_k)-F(x^\star)\leq\frac{\|x_0-x^\star\|_2^2}{2tk},
 \qquad k\geq1.
\]
\end{proposition}
\begin{proof}
记 $x^+$ 为 $Q_x$ 的最小点。由 $Q_x$ 的 $1/t$-强凸性，对任何 $y\in\operatorname{dom}R$ 有
\[
 Q_x(x^+)+\frac1{2t}\|y-x^+\|_2^2\leq Q_x(y).
\]
此式也可不使用次微分直接证明：对 $x^+$ 与 $y$ 的凸组合应用强凸不等式，再用 $Q_x(x^+)$ 不大于组合点的值，最后令组合中 $y$ 的权重趋于零。
又有 $F(x^+)\leq Q_x(x^+)$，以及由 $h$ 的凸性得到的
$Q_x(y)\leq F(y)+\|y-x\|_2^2/(2t)$。相合便得到与式\eqref{eq:co-projected-three-point}相同的三点不等式，其中 $f$ 换为 $F$。取 $y=x$ 得目标不增，取 $y=x^\star$ 后累加即得结果。
\end{proof}

对 Lasso，$h(x)=\|Ax-b\|_2^2/2$，当 $A\ne0$ 时可取 $L=\|A\|_2^2$。算法从给定 $x_0$ 出发，每步计算残差、梯度 $A^{\mathsf T}(Ax_k-b)$，再作一次软阈值。$A=0$ 的退化情形可直接最小化正则项，无需采用 $1/L$。例如 $A=I$、$b=(3,1/2)$、$\alpha=1$、$t=1$，任意初值的试探点都为 $b$，一次软阈值就到达 $(2,0)$。

一般情况下可以计算近端梯度映射
\[
 \mathcal G_t(x)=\frac{x-\operatorname{prox}_{tR}(x-t\nabla h(x))}{t}.
\]
近端子问题的一阶条件给出 $\mathcal G_t(x)=0$ 当且仅当
$0\in\nabla h(x)+\partial R(x)$，也就是复合问题的最优性条件。实际算法应给出容差与最大迭代数；小映射范数的退出表示近似满足这一条件，若要声明目标误差不超过指定值，还应给出相应的误差界。

\begin{example}[给 Lasso 迭代附上可计算的停止证书]
设 $\alpha>0$，当前点为 $x$，原目标为 $F(x)=\|Ax-b\|_2^2/2+\alpha\|x\|_1$。构造一个不需要预先知道 $x^\star$ 的目标误差上界。
\end{example}
\begin{solve}
引入残差变量 $r=Ax-b$，给等式 $Ax-b-r=0$ 配乘子 $u$。拉格朗日函数为
\[
 \frac12\|r\|_2^2+\alpha\|x\|_1+u^{\mathsf T}(Ax-b-r).
\]
对 $r$ 取下确界得到 $-\|u\|_2^2/2$；对 $x$ 取下确界在 $\|A^{\mathsf T}u\|_\infty\leq\alpha$ 时为零，否则为 $-\infty$。所以对偶问题为
\[
 \max_u d(u)=-\frac12\|u\|_2^2-b^{\mathsf T}u,
 \qquad \|A^{\mathsf T}u\|_\infty\leq\alpha.
\]
计算当前残差 $r=Ax-b$，令
\[
 c=\max\{1,\|A^{\mathsf T}r\|_\infty/\alpha\},\qquad u=r/c.
\]
该 $u$ 必对偶可行，弱对偶直接给出 $0\leq F(x)-F^\star\leq F(x)-d(u)$。因此近端梯度迭代可以在这个可计算的差不超过 $\varepsilon$ 时停止。对 $A=I,b=(3,1/2),\alpha=1,x=(2,0)$，有 $u=(-1,-1/2)$，原对偶值同为 $21/8$，间隙为零，精确验证了前面的答案。
\end{solve}

投影、软阈值和一般近端算子的系统讨论可参见\cite{parikh-boyd}。本节只使用其中最基本的模型与精确迭代；若内层近端点近似计算，需把内层误差纳入收敛分析，不能直接照用上面的精确界。

\section{次梯度法与非光滑性的代价}
\label{sec:co-subgradient-method}

当整个目标都非光滑，或者近端子问题不易求解时，可以使用次梯度法（subgradient method）。设 $f:\mathbb R^n\to\mathbb R$ 凸，每步选择 $g_k\in\partial f(x_k)$，按
\[
 x_{k+1}=x_k-t_k g_k,\qquad t_k>0
\]
更新。次梯度是全局仿射下界的斜率，负次梯度却不一定是下降方向。例如 $f(x)=|x|$ 在零点允许取 $g=1$，沿负次梯度移动反而使目标增大；在该点也允许选到零次梯度以识别最优性。因此一般不能把光滑函数的下降线搜索直接移用过来。

\begin{proposition}[有界次梯度下的最佳历史值界]
设最优解 $x^\star$ 存在，$\|x_0-x^\star\|_2\leq R$，且所选次梯度满足 $\|g_k\|_2\leq G$。对任意 $N\geq1$，
\[
 \min_{0\leq k<N}\bigl(f(x_k)-f^\star\bigr)
 \leq\frac{R^2+G^2\sum_{k=0}^{N-1}t_k^2}
              {2\sum_{k=0}^{N-1}t_k}.
\]
\end{proposition}
\begin{proof}
展开点误差平方，并用次梯度不等式
$g_k^{\mathsf T}(x_k-x^\star)\geq f(x_k)-f^\star$，得到
\[
 \|x_{k+1}-x^\star\|_2^2
 \leq\|x_k-x^\star\|_2^2
 -2t_k(f(x_k)-f^\star)+t_k^2G^2.
\]
求和后舍去左端非负项，再用各项误差都不小于历史最小误差，便得结论。
\end{proof}

当 $R,G>0$ 且预先给定迭代总数 $N$，取恒定步长 $t_k=R/(G\sqrt N)$，便得到误差界 $RG/\sqrt N$。这个步长依赖计划运行的总步数；如果无限运行却一直保持某个固定正步长，上述界一般不会趋于零。另一种选择是让 $\sum_k t_k=+\infty$、$\sum_k t_k^2<+\infty$，从而历史最佳目标值趋于最优值。此处没有断言每一步目标都下降，也没有仅由该不等式证明全部迭代点收敛。

算法应保存目标值最低的历史点，或输出平均点 $\bar x_N=(\sum_{k=0}^{N-1}t_kx_k)/(\sum_{k=0}^{N-1}t_k)$；凸性保证后者也满足同一上界。给定有限预算时，到达最大步数后可报告已有理论界；未知 $R,G$ 时，则不能把上述具体精度当成已验证的数值保证。与光滑梯度法的 $O(1/k)$ 比较，这里的 $O(1/\sqrt{k})$ 展示了在这些假设下非光滑一阶信息所能提供的较弱基本保证。

\section{最速下降、尺度与条件数}

给定范数 $\|\cdot\|$，归一化最速下降方向（steepest descent direction）最小化 $g^{\mathsf T}d$，约束 $\|d\|\leq1$。其最小值为 $-\|g\|_*$，其中对偶范数为 $\|g\|_*=\max_{\|d\|\leq1}g^{\mathsf T}d$。当 $g\ne0$ 时，欧氏范数下的归一化方向为 $-g/\|g\|_2$，与负梯度平行；$\ell_1$ 范数对应沿最大绝对梯度分量的坐标轴移动。

在二次范数 $\|d\|_P=\sqrt{d^{\mathsf T}Pd}$、$P\succ0$ 下，方向与 $-P^{-1}g$ 平行。矩阵 $P$ 改变所用度量，称为预条件（preconditioning）的一种形式；它不改变原目标的最优点。

对正定二次函数，梯度的变化常数可取最大特征值 $L$，强凸参数可取最小特征值 $\mu$，比值 $\kappa=L/\mu$ 为谱条件数。例 $f(x)=\tfrac12(x_1^2+100x_2^2)$，步长 $1/100$ 使第二个分量一步归零，却使第一个分量每步只乘以 $0.99$。选择 $P=\operatorname{diag}(1,100)$ 后，预条件方向配单位步长可一步到达最优点。

\section{牛顿方向与阻尼牛顿法}

设 $f$ 二次连续可微，当前 $H=\nabla^2f(x)\succ0$。二阶模型为 $m(d)=f(x)+g^{\mathsf T}d+d^{\mathsf T}Hd/2$，其唯一最小点满足
\begin{equation}
 Hd=-g.
 \label{eq:co-newton-direction}
\end{equation}
这定义牛顿方向（Newton direction）。当 $g\ne0$ 时，$g^{\mathsf T}d=-g^{\mathsf T}H^{-1}g<0$，故可用回溯线搜索形成阻尼牛顿法（damped Newton method）。实现应解线性系统，不显式计算逆矩阵。

\begin{example}[牛顿步为何需要阻尼]
在 $x>0$ 上最小化 $f(x)=x-\log x$，其唯一最优点为 $x^\star=1$。
\end{example}
\begin{solve}
梯度为 $1-1/x$，海塞矩阵为 $1/x^2$，所以牛顿方向为 $d=x-x^2$，单位步长更新为 $x^+=2x-x^2$。若 $x=3$，更新得到 $-3$，已经离开定义域；回溯必须先要求 $3-6t>0$，再检查充分下降。在 $x=1+e$ 且 $e$ 很小时，更新误差为 $e^+=-e^2$，表现为误差平方级缩小。相同的牛顿公式在远处可能不可行，在近处却很快，这正是全局阻尼与局部二次收敛需要分开讨论的原因。
\end{solve}

\begin{theorem}[局部二次收敛]
设 $x^\star$ 为驻点，在其某邻域内 $\nabla^2f(x)\succeq\mu I$、$\mu>0$，且海塞矩阵关于谱范数以常数 $M$ 利普希茨连续。对充分接近 $x^\star$ 的初值，单位步长牛顿法保持在该邻域内，并满足
\[
 \|x_{k+1}-x^\star\|_2\leq\frac{M}{2\mu}\|x_k-x^\star\|_2^2.
\]
\end{theorem}
\begin{proof}
记 $e=x-x^\star$、$H=\nabla^2f(x)$。因 $\nabla f(x^\star)=0$，积分表示给出
\[
 e^+=H^{-1}\int_0^1[H-\nabla^2f(x^\star+te)]e\,dt.
\]
用 $\|H^{-1}\|_2\leq1/\mu$ 及海塞矩阵的利普希茨界，得到 $\|e^+\|_2\leq(M/(2\mu))\|e\|_2^2$。选初值足够近，使此上界既留在邻域内又小于原误差，即可归纳。
\end{proof}

局部定理不保证任意初值都适合单位步长。若初值次水平集紧、包含在定义域内，并在其邻域内有统一界 $0<\mu I\preceq\nabla^2f\preceq LI$，则回溯接受步长有正的统一下界，每次下降量可下界为常数乘 $\|\nabla f\|_2^2$，由目标有下界可推出梯度趋零；再由强凸性推出收敛到唯一解。正定二次目标的二阶模型本身精确，解准牛顿方程并取单位步长即可一步求解。

牛顿法可按以下流程实现：给定定义域内的 $x_0$、线搜索参数及容差；在每步计算梯度和海塞矩阵，检查是否满足所用曲率条件，解 $H_kd_k=-g_k$；若已有适用的误差证书则检查停止，否则用回溯选取保持定义域和充分下降的步长；更新后重新计算。若 $H_k$ 奇异或非正定，原牛顿方向的定义和下降性证明不再适用。加入 $\rho I$ 的正则化牛顿步是一种不同的局部模型，需要说明 $\rho$ 的选择，不能把修正后的方向仍当作原方程的精确解。

\section{牛顿减量与自协调分析}

本节为选读。此前的分析用全局常数 $L,\mu$ 控制曲率，而障碍函数在边界附近曲率可能无界，找一个整个定义域都适用的 $L$ 很困难。自协调分析改为使用当前位置的海塞矩阵来衡量位移，再控制沿这段位移曲率能变化多少。这是它适合内点法的主要原因。

在 $H\succ0$ 时定义牛顿减量（Newton decrement）
\[
 \delta(x)=\sqrt{\nabla f(x)^{\mathsf T}H^{-1}\nabla f(x)}.
\]
二阶模型预测的最大下降量为 $\delta(x)^2/2$，因为将 $d=-H^{-1}g$ 代入 $m(d)$ 可得此值。但它本身不是任意函数的全局目标误差界。

自协调函数（self-concordant function，亦译自和谐函数）是开凸域上的三次连续可微凸函数，满足
\[
 |D^3f(x)[h,h,h]|\leq2\bigl(D^2f(x)[h,h]\bigr)^{3/2}
\]
对所有点和方向成立，其中 $D^j f(x)[h,\ldots,h]$ 为沿直线的 $j$ 阶导数。$-\log x$ 在正半轴上恰好满足等号；凸二次函数的三阶导数为零。该条件用局部曲率控制高阶变化，且保持仿射坐标变换下的形式。

以下标准自协调估计作为选读结论引用，完整推导见\cite{boyd-vandenberghe}关于牛顿法的讨论：若 $f$ 为自协调函数，在域外取 $+\infty$ 后闭、海塞矩阵正定且最小值取到，则步长 $1/(1+\delta)$ 给出至少 $\delta-\log(1+\delta)$ 的目标下降；当 $\delta<1$ 时，
\[
 f(x)-f^\star\leq-\delta-\log(1-\delta).
\]
因此在这些条件下，减量可以转成严格误差证书。自协调性也为内点法的迭代分析提供工具。

\section{实现、停止准则与数值实践}

梯度法每步需要一个梯度和若干函数值；牛顿法还需要海塞矩阵及线性求解。稠密 $n\times n$ 正定矩阵的 Cholesky 分解需 $O(n^3)$ 运算、$O(n^2)$ 存储；稀疏结构、低秩结构和分解复用可能显著改变成本。迭代次数少不必然意味着总时间短。

在已知强凸参数 $\mu$ 时，可用 $\|\nabla f(x)\|_2^2/(2\mu)\leq\varepsilon$ 停止；适用的自协调情形可用前节减量界。缺少误差界时，小梯度只直接表示近似驻点。实际实现还需设置最大迭代次数与线搜索失败状态，不能把这些退出状态等同于已满足最优性精度。

\begin{enumerate}
 \item 推导 $f(x)=x^2/2$ 上固定步长梯度法的序列。证明 $0<t<2$ 时收敛，$t=2$ 时一般振荡。
 \item 对 $f(x)=\tfrac12(x_1^2+100x_2^2)$，比较固定步长梯度法与牛顿法，记录目标误差和线性求解代价。
 \item 验证 $-\log x$ 的自协调条件，解释它虽然自协调却没有有限最小点，故不能直接套用含最小点的误差结论。
 \item 证明若 $\mu$ 已知且目标误差不超过 $\varepsilon$，则点误差不超过 $\sqrt{2\varepsilon/\mu}$。
 \item 基础练习：对 $\min_{x\geq0}(x+1)^2/2$ 从 $x_0=2$ 作步长 $t=1/2$ 的投影梯度迭代，列出前两步。答案为 $x_1=1/2,x_2=0$，此后保持为零。
 \item 推导练习：对 $F(x)=(x-3)^2/2+|x|$，从 $x_0=0$ 作步长 $t=1/2$ 的近端梯度迭代。求得 $x_1=1,x_2=3/2$，证明 $x_k\to2$，并计算一个合法对偶点及其间隙。
 \item 比较练习：对 $|x|$ 使用固定步长次梯度法，从 $x_0=1/4$、$t=1/2$ 开始，非零处取符号次梯度。证明它在 $\pm1/4$ 间振荡，解释为什么这不违反本章关于递减步长或有限预算的结论。
\end{enumerate}

\chapter{等式约束优化算法}
\label{ch:co-equality}

考虑 $\min_x f(x)$，满足 $Ax=b$。等式要求每一步位移保持在可行方向内，或在迭代中逐步消除等式残差。本章假设 $f$ 在开凸域上二次连续可微，并在使用牛顿系统时明确秩与曲率条件。

例如要求 $x_1+x_2=1$ 时，可行位移必须满足 $d_1+d_2=0$。即使目标沿横轴下降，也不能只改变第一分量；必须同时补偿第二分量。等式约束算法的任务，是从一开始就在这样的可行方向中寻找改善，或者同时减少目标驻点残差和等式误差。

\section{仿射约束与降维表示}

设 $A\in\mathbb R^{p\times n}$。首先检查 $b\in\operatorname{range}A$；否则等式系统不可行。可以删除冗余等式，以下通常假设 $A$ 满行秩。若 $p=n$，可行点唯一，只需检查其是否在目标定义域内。

当 $p<n$，取可行点 $\bar x$，令 $F\in\mathbb R^{n\times(n-p)}$ 的列构成 $\ker A$ 的一组基。所有可行点唯一写成 $x=\bar x+Fz$，于是原问题化为无约束问题
\[
 \min_z\widetilde f(z)=f(\bar x+Fz),\qquad
 \nabla\widetilde f=F^{\mathsf T}\nabla f,\quad
 \nabla^2\widetilde f=F^{\mathsf T}\nabla^2fF.
\]
因此只有可行方向上的曲率直接影响约束最优性。

在约束 $x_1+x_2=1$ 上，可以取 $\bar x=(1,0)^{\mathsf T}$、$F=(1,-1)^{\mathsf T}$，于是 $x=(1+z,-z)^{\mathsf T}$。二维问题实际只有一个自由变量。等式约束并不一定增加数学难度，有时恰好能消去多余的自由度。例 $f(x_1,x_2)=x_1^2/2$、$x_2=0$，全空间海塞矩阵奇异，但可行子空间上的目标严格凸。

对于凸可微目标，约束最优性等价于 $F^{\mathsf T}\nabla f(x^\star)=0$。这是说梯度与所有可行方向正交，而不是说梯度本身为零。由于 $(\ker A)^\perp=\operatorname{range}A^{\mathsf T}$，它又等价于存在 $\nu$ 使 $\nabla f(x^\star)+A^{\mathsf T}\nu=0$。零空间形式和拉格朗日乘子形式只是同一几何事实的两种坐标表达。

\section{等式约束二次问题与 KKT 系统}

对二次目标 $f(x)=x^{\mathsf T}Px/2+q^{\mathsf T}x+r$、$P=P^{\mathsf T}$，驻点与可行性构成
\[
 \begin{bmatrix}P&A^{\mathsf T}\\A&0\end{bmatrix}
 \begin{bmatrix}x\\\nu\end{bmatrix}
 =\begin{bmatrix}-q\\b\end{bmatrix}.
\]
即使 $P$ 在全空间不正定，只要在 $\ker A$ 上正定，该二次问题也有唯一约束最小点。

\begin{proposition}[KKT 矩阵的非奇异性]
若 $A$ 满行秩，且对所有非零 $d\in\ker A$ 有 $d^{\mathsf T}Pd>0$，则上述 KKT 矩阵非奇异。
\end{proposition}
\begin{proof}
设齐次系统满足 $Pd+A^{\mathsf T}w=0$、$Ad=0$。第一式左乘 $d^{\mathsf T}$ 得 $d^{\mathsf T}Pd=0$，故 $d=0$；随后 $A^{\mathsf T}w=0$，由满行秩得 $w=0$。齐次系统仅有零解，因此矩阵可逆。
\end{proof}
冗余约束会使乘子不唯一，即使原变量的最优点仍唯一。因此数值奇异可能来自约束表示，并不必然来自目标缺少约束最小值。

\section{可行初始点的约束牛顿法}

给定 $Ax_k=b$，令 $g_k=\nabla f(x_k)$、$H_k=\nabla^2f(x_k)$。最小化局部二次模型，约束方向满足 $Ad=0$，得到
\begin{equation}
 \begin{bmatrix}H_k&A^{\mathsf T}\\A&0\end{bmatrix}
 \begin{bmatrix}d_k\\w_k\end{bmatrix}
 =-\begin{bmatrix}g_k\\0\end{bmatrix}.
 \label{eq:co-equality-newton}
\end{equation}
若 $H_k$ 在 $\ker A$ 上正定，系统唯一可解。由方向方程左乘 $d_k^{\mathsf T}$ 得 $g_k^{\mathsf T}d_k=-d_k^{\mathsf T}H_kd_k\leq0$；非零方向严格下降，零方向对应约束驻点。

算法从可行初值开始，每步解式\eqref{eq:co-equality-newton}，用回溯保证目标下降及定义域可行，然后更新 $x_{k+1}=x_k+t_kd_k$。由于 $Ad_k=0$，任意步长均保持等式。约束牛顿减量可定义为 $\sqrt{d_k^{\mathsf T}H_kd_k}$，其为零当且仅当约束方向的一阶最优性条件成立；转换为误差上界仍需曲率或自协调假设。

令 $d_k=F\Delta z_k$，将方向方程左乘 $F^{\mathsf T}$，得到 $F^{\mathsf T}H_kF\Delta z_k=-F^{\mathsf T}g_k$，这就是降维目标的牛顿方程。因此第\ref{ch:co-unconstrained}章的局部收敛定理可在降维空间中使用，而无需全空间海塞矩阵正定。

\section{不可行初始点与原对偶残差}

若初值不满足等式，令原残差（primal residual）和对偶残差（dual residual）为
\[
 r_{\mathrm p}(x)=Ax-b,\qquad
 r_{\mathrm d}(x,\nu)=\nabla f(x)+A^{\mathsf T}\nu.
\]
对二者组成的非线性方程使用牛顿法，得到
\[
 \begin{bmatrix}\nabla^2f(x)&A^{\mathsf T}\\A&0\end{bmatrix}
 \begin{bmatrix}\Delta x\\\Delta\nu\end{bmatrix}
 =-\begin{bmatrix}r_{\mathrm d}\\r_{\mathrm p}\end{bmatrix}.
\]
更新 $(x,\nu)\leftarrow(x,\nu)+t(\Delta x,\Delta\nu)$。原残差精确满足 $r_{\mathrm p}^{+}=(1-t)r_{\mathrm p}$；若取单位步长，下一点即满足等式，但可能因目标定义域或非线性残差而不能采用该步长。

可对总残差 $R=(r_{\mathrm d},r_{\mathrm p})$ 回溯，要求 $\|R^+\|_2\leq(1-\alpha t)\|R\|_2$ 并保持定义域。若残差雅可比可逆且方向精确，则 $R^+=R-tR+o(t)$，故小步长会被接受。这里下降的是残差，不保证目标每步都下降。局部快速收敛要求雅可比非奇异和相应光滑性，全局收敛仍需额外条件。

\begin{example}[从不可行点同时修正两个残差]
仍取 $f(x)=[(x_1-2)^2+x_2^2]/2$、$x_1+x_2=1$，但从 $x=(0,0)$、$\nu=0$ 开始。
\end{example}
\begin{solve}
当前 $r_{\mathrm p}=-1$、$r_{\mathrm d}=(-2,0)^{\mathsf T}$。牛顿方程为
\[
 \Delta x_1+\Delta\nu=2,\qquad
 \Delta x_2+\Delta\nu=0,\qquad
 \Delta x_1+\Delta x_2=1.
\]
解得 $\Delta\nu=1/2$、$\Delta x=(3/2,-1/2)$。采用单位步长后，$x=(3/2,-1/2)$、$\nu=1/2$，原残差和对偶残差都为零。这里残差方程本身是仿射的，故牛顿线性化没有近似误差；对一般非二次目标，单位步长只保证等式残差消失，不保证驻点残差同时消失。
\end{solve}

\section{块消元与结构化求解}

若 $H\succ0$，方向系统 $H\Delta x+A^{\mathsf T}\Delta\nu=-r_{\mathrm d}$、$A\Delta x=-r_{\mathrm p}$ 可消去 $\Delta x$，得到
\[
 AH^{-1}A^{\mathsf T}\Delta\nu
 =r_{\mathrm p}-AH^{-1}r_{\mathrm d},\qquad
 \Delta x=-H^{-1}(r_{\mathrm d}+A^{\mathsf T}\Delta\nu).
\]
满行秩使 Schur 补 $AH^{-1}A^{\mathsf T}$ 正定。公式中的 $H^{-1}$ 在实现时由解方程完成。若 $H$ 仅在零空间上正定，可能不可逆，应使用零空间法或对整个对称不定 KKT 矩阵作分解。

零空间法降低变量维数，但构造稠密基可能破坏原稀疏结构；Schur 补法适合等式较少的情形，却可能增加填充和条件数。实际方法的选择取决于维数与矩阵结构，不能只比较公式长短。

\section{例题、终止条件与学习检查}
\begin{example}[守恒约束下的一步牛顿]
最小化 $f(x)=\tfrac12[(x_1-2)^2+x_2^2]$，满足 $x_1+x_2=1$。取可行初值 $x_0=(1,0)$，有 $g_0=(-1,0)$、$H_0=I$。方向方程为 $d_1+w=1$、$d_2+w=0$、$d_1+d_2=0$，得 $w=1/2$、$d=(1/2,-1/2)$。单位步长到达 $x^\star=(3/2,-1/2)$，目标值为 $1/4$；这是二次模型精确性在约束空间中的表现。
\end{example}

停止时至少检查 $\|Ax-b\|_2$ 与 $\|\nabla f(x)+A^{\mathsf T}\nu\|_2$。尺度不同的残差应结合绝对、相对容差解释；当 $x$ 已可行且 $f$ 在全空间为 $\mu$-强凸时，$\|r_{\mathrm d}\|_2^2/(2\mu)$ 是目标误差上界，因为 $A(x-x^\star)=0$ 使乘子项消去。不可行点不能直接使用这个结论。

\begin{enumerate}
 \item 用消元 $x_2=1-x_1$ 重新求解上述例题，并核对方向。
 \item 在 $f(x)=x_1^2/2$、$x_2=0$ 中写出 KKT 矩阵，验证其可逆而 $\nabla^2f$ 不可逆。
 \item 给等式系统重复添加一条相同约束，说明原解、乘子和线性系统非奇异性分别怎样变化。
 \item 对不可行初始点方法验证原残差递推，并说明单位步长后的目标域检查为什么仍然必要。
 \item 推导练习：若 $F$ 的列正交归一，说明 $\|F^{\mathsf T}\nabla f(x)\|_2$ 衡量的是哪个可行子空间中的梯度；在 $x_1+x_2=1$ 的二维例中直接计算它，并与乘子残差比较。
\end{enumerate}

\chapter{内点法}
\label{ch:co-interior}

内点法（interior-point method）把不等式约束转化为一系列内部子问题，再用牛顿法求解。对偶理论给出这些内部点距原最优值的界，因而将算法更新与精度认证联系起来。

\section{不等式约束与严格可行域}

考虑 $\min f_0(x)$，满足 $f_i(x)\leq0$、$Ax=b$，本章设不等式数量 $m\geq1$，$f_0,\ldots,f_m$ 在公共开凸域上二次连续可微且凸，$A$ 满行秩。严格可行域为满足全部 $f_i(x)<0$ 和仿射等式的点。基本障碍法从这样的点出发；没有严格可行点时，应先消去冗余约束、降低维数或采用其他问题表示。

原最优点可能位于边界，例如 $\min_{x\geq0}x$ 的解为零。内点法在内部逐步逼近它，不要求某个有限迭代恰好到达边界。算法能否达到指定目标误差，还依赖子问题的存在性与求解精度。

内点法与投影法处理边界的方式不同。投影法可以先走到外部再投影到边界；障碍法则通过一个在边界附近急剧增加的代价，使内部迭代避免碰到边界。这里“内部”是相对于仿射等式所确定的可行空间而言，不要求等式可行点拥有全空间中的开邻域。

\section{对数障碍与中心路径}

在严格不等式域上定义对数障碍（logarithmic barrier）
\[
 \phi(x)=-\sum_{i=1}^m\log(-f_i(x)).
\]

量 $s_i=-f_i(x)>0$ 是第 $i$ 条不等式的余量。$-\log s_i$ 在 $s_i\downarrow0$ 时趋于 $+\infty$，于是过分贴近任何一条边界都会付出很大代价。对于 $x\geq0$，障碍就是 $-\log x$；对于 $0\leq x\leq1$，障碍为 $-\log x-\log(1-x)$，同时阻止越过两个端点。
$-\log(-u)$ 在 $u<0$ 上凸且递增，所以它与凸 $f_i$ 的复合凸。直接求导得到
\[
 \nabla\phi=-\sum_i\frac{\nabla f_i}{f_i},\qquad
 \nabla^2\phi=\sum_i\left(\frac{\nabla f_i\nabla f_i^{\mathsf T}}{f_i^2}
                         -\frac{\nabla^2f_i}{f_i}\right)\succeq0.
\]
对 $t>0$，中心化子问题为
\begin{equation}
 x(t)\in\mathop{\mathrm{arg\,min}}_{Ax=b}\{t f_0(x)+\phi(x)\}.
 \label{eq:co-central-path}
\end{equation}

除以 $t$ 后，同一个问题写成最小化 $f_0(x)+\phi(x)/t$。所以增大 $t$ 等价于逐渐减弱障碍的权重，而不是增强它。较小的 $t$ 更重视与边界保持距离，较大的 $t$ 更重视原目标；极限中才可能逼近位于边界的原最优点。
假设这些解存在；若目标在可行子空间上严格凸，则解唯一，由这些点构成中心路径（central path）。

\begin{example}[一维中心路径]
对 $\min_{x\geq0}x$，障碍子问题为 $\min_{x>0}tx-\log x$。驻点给出 $x(t)=1/t$，对偶乘子恒为 $1$，最优值下界为零，间隙恰为 $1/t$。内部点随 $t\to\infty$ 逼近边界解零。
\end{example}

\begin{theorem}[精确中心点的误差界]
若 $x(t)$ 是精确中心点，则可构造对偶可行乘子，使原对偶间隙为 $m/t$，从而 $0\leq f_0(x(t))-p^\star\leq m/t$。
\end{theorem}
\begin{proof}
中心化驻点条件为 $t\nabla f_0(x(t))+\nabla\phi(x(t))+A^{\mathsf T}w=0$。令
\[
 \lambda_i(t)=-\frac1{t f_i(x(t))}>0,\qquad\nu(t)=w/t.
\]
则原拉格朗日函数在 $x(t)$ 处梯度为零，由凸性知该点取到其下确界。因此
\[
 g(\lambda(t),\nu(t))=f_0(x(t))+\sum_i\lambda_i(t)f_i(x(t))
 =f_0(x(t))-m/t.
\]
弱对偶即给出目标误差上界。
\end{proof}
这同时解释了互补松弛的近似形式 $\lambda_i(t)(-f_i(x(t)))=1/t$。若中心点只被近似求出，拉格朗日最小化不再精确成立，不能不计内层误差便断言同样的间隙等式。

\begin{example}[区间上的中心点与误差证书]
考虑 $\min x$，满足 $0\leq x\leq1$。求其中心路径，并解释 $m/t$ 代表什么。
\end{example}
\begin{solve}
中心化目标为 $tx-\log x-\log(1-x)$，导数为 $t-1/x+1/(1-x)$，海塞矩阵恒正。解驻点式得到
\[
 x(t)=\frac{t+2-\sqrt{t^2+4}}{2t}
     =\frac{2}{t+2+\sqrt{t^2+4}}\in(0,1/2).
\]
后一形式避免大数相减，数值计算更稳定。$t\downarrow0$ 时它趋于区间中点 $1/2$；$t\to+\infty$ 时趋于原最优点零。两条不等式的乘子分别为
\[
 \lambda_1=\frac1{tx(t)},\qquad
 \lambda_2=\frac1{t(1-x(t))}.
\]
它们满足 $1-\lambda_1+\lambda_2=0$，故原对偶间隙为 $2/t$。但原目标误差是 $x(t)$，通常严格小于 $2/t$；证书给的是上界，不必等于真实误差。在 $t=10$ 时，$x(t)\approx0.0901$，证书为 $0.2$。
\end{solve}

若只有近似中心点，也可以在额外曲率条件下补偿内层误差。设当前 $x$ 严格可行且 $Ax=b$，取 $\lambda_i=-1/(tf_i(x))$ 和任意 $\nu$。假设对应的 $\mathcal L(\cdot,\lambda,\nu)$ 在公共定义域上为已知 $\mu_{\mathcal L}>0$ 的强凸函数，记 $r=\nabla_x\mathcal L(x,\lambda,\nu)$。强凸下界给出
\[
 g(\lambda,\nu)\geq
 \mathcal L(x,\lambda,\nu)-\frac{\|r\|_2^2}{2\mu_{\mathcal L}},
\]
因此
\begin{equation}
 f_0(x)-p^\star
 \leq\frac mt+\frac{\|r\|_2^2}{2\mu_{\mathcal L}}.
 \label{eq:co-inexact-center-bound}
\end{equation}
右侧把中心路径误差与内层驻点误差分开。这个界要求明确的强凸常数；一般模型若没有该常数，应使用可直接计算的对偶函数或其他适用的中心化估计。

\section{障碍方法的两层迭代}

给定严格可行点 $x_0$、参数 $t_0>0$、增长因子 $\gamma>1$ 和最终精度 $\varepsilon>0$，基本障碍法如下。
\begin{enumerate}
 \item 在当前 $t$ 下，用第\ref{ch:co-equality}章的约束牛顿法最小化 $tf_0+\phi$。线搜索同时保持所有 $f_i(x)<0$。
 \item 检查适用的中心化误差界及原问题误差证书。精确中心化的理想模型在 $m/t\leq\varepsilon$ 时可停止。
 \item 否则令 $t\leftarrow\gamma t$，以上一个中心点热启动（warm start）新的子问题，重复迭代。
\end{enumerate}
理想精确中心化下，参数更新次数至多为 $\max\{0,\lceil\log(m/(t_0\varepsilon))/\log\gamma\rceil\}$。这只是外层次数；每层还需要若干牛顿步和线性求解。增长因子大减少外层次数，却可能让新子问题的初值离中心更远。

\section{阶段一问题与可行性判断}

阶段一方法（phase I method）通过辅助变量寻找严格可行点。例如求解
\[
 \min_{x,s}s,\qquad f_i(x)\leq s\ (i=1,\ldots,m),\quad Ax=b.
\]
若已经有定义域内的等式可行点 $x_0$，选择 $s_0>\max_i f_i(x_0)$ 即得到辅助问题的严格可行初值。任意目标值 $s<0$ 的辅助可行点都给出原问题严格可行点，无须继续求到辅助最优值。

若辅助问题的最优值严格为正，则原问题不可行，因为原可行点必可与 $s=0$ 配成辅助可行点。若最优值为零且取到，则原问题可行，但没有严格可行点；若零下确界未取到，则仅凭数值接近零不能判断原可行性。实际计算需要适用的对偶证书来区分不可行与尚未找到可行点。

\begin{example}[阶段一中零与正数的区别]
比较系统 $x\leq0,\ -x\leq0$ 与系统 $x\leq0,\ 1-x\leq0$。
\end{example}
\begin{solve}
第一个系统只有可行点零。其阶段一目标为 $\min_x\max\{x,-x\}=0$，在零点取到，但不存在使两条不等式同时严格成立的点。它可以等价改写为等式 $x=0$。第二个系统要求同时 $x\leq0$ 和 $x\geq1$，不可行；其阶段一目标 $\max\{x,1-x\}$ 的最小值为 $1/2>0$，在 $x=1/2$ 取到。两个问题都无法找到原始表示下的严格可行点，数学原因却不同。
\end{solve}

若只算得辅助可行点 $s=10^{-6}$，它只说明当前最大违约程度很小，没有证明辅助最优值为正。要认证原问题不可行，需要一个严格为正的辅助最优值下界；这又回到了对偶证书的作用。

\section{自协调障碍与复杂度}

本节为选读，目标是解释复杂度结论依赖的条件。掌握基本障碍法及误差证书后，可以暂时跳到原对偶方程；不要把这里的迭代次数界理解成任何凸模型在任何初始化下都具有相同运行时间。

对开凸集 $C$，自协调障碍除满足第\ref{ch:co-unconstrained}章的三阶导数条件外，还在趋近边界时发散，并具有某个参数 $\vartheta>0$，使
\[
 |D\phi(x)[h]|\leq\sqrt{\vartheta}\sqrt{D^2\phi(x)[h,h]}
\]
对所有点和方向成立。参数控制障碍梯度相对于局部曲率的大小。线性目标 $t c^{\mathsf T}x+\phi(x)$ 的三阶与二阶导数不受线性项影响，所以可统一应用自协调牛顿分析。

一个标准短步路径跟踪（short-step path-following）结论如下：在有界开凸可行域上，给定海塞矩阵正定的 $\vartheta$-自协调障碍，消去仿射等式后从 $t_0$ 的中心点或指定中心邻域出发，每次将 $t$ 乘以 $1+\eta/\sqrt\vartheta$，其中 $\eta>0$ 为足够小的固定常数，并用牛顿步维持固定的中心邻域，则达到原对偶间隙不超过 $\varepsilon$ 所需牛顿步数为
\[
 O\!\left(\sqrt\vartheta\,
 \log\max\left\{2,\frac{\vartheta}{t_0\varepsilon}\right\}\right).
\]
此为选读复杂度结论，完整中心邻域估计见\cite{boyd-vandenberghe}的自协调分析。证明要点是每次参数变化只使中心偏移一个受控的局部距离，因而常数次牛顿校正足够；达到 $t$ 与 $\vartheta/\varepsilon$ 同阶所需更新次数即产生上式。

该界不包含取得初始中心点的代价，也不等于浮点总运算量。一般凸不等式的对数障碍未必满足所需自协调条件；上一节任意固定大增长因子的障碍法也不能直接套用短步法的迭代界。

\section{锥约束的障碍函数}

三种常用锥的标准障碍及参数为：
\[
 \begin{array}{c|c|c}
 \text{定义域}&\text{障碍函数}&\vartheta\\\hline
 x\in\mathbb R_{++}^n&-\sum_i\log x_i&n\\
 t>\|x\|_2&-\log(t^2-\|x\|_2^2)&2\\
 X\in\mathbb S_{++}^n&-\log\det X&n
 \end{array}
\]
这些是相应锥的自协调障碍，矩阵情形采用迹内积；这里引用其标准性质。积锥的障碍相加，参数相加。上述函数还满足 $\phi(\alpha x)=\phi(x)-\vartheta\log\alpha$，$\alpha>0$，称为对数齐次性（logarithmic homogeneity）。对 $\alpha$ 求导于 $1$，得到 $\langle\nabla\phi(x),x\rangle=-\vartheta$。

对于变量位于锥内的线性目标问题，精确中心点对应的对偶锥变量为 $z=-\nabla\phi(x)/t$，其互补内积为 $\langle z,x\rangle=\vartheta/t$。因此锥形式的间隙由障碍参数控制，不能一律用标量不等式的条数代替。

\section{原对偶内点法}

引入松弛向量 $s>0$，使 $f(x)+s=0$，并保持乘子 $\lambda>0$。这里 $f=(f_1,\ldots,f_m)$，$J(x)$ 为以 $\nabla f_i(x)^{\mathsf T}$ 作行的雅可比矩阵，$S=\operatorname{diag}(s)$、$\Lambda=\operatorname{diag}(\lambda)$。定义
\[
 \begin{aligned}
 r_{\mathrm d}&=\nabla f_0+J^{\mathsf T}\lambda+A^{\mathsf T}\nu,&
 r_{\mathrm e}&=Ax-b,\\
 r_{\mathrm i}&=f(x)+s,&
 r_{\mathrm c}&=S\lambda-\sigma\bar\mu\mathbf1,
 \end{aligned}
 \qquad \bar\mu=\frac{s^{\mathsf T}\lambda}{m},\quad0<\sigma<1.
\]
将本步目标 $\sigma\bar\mu$ 固定，设 $H=\nabla^2f_0+\sum_i\lambda_i\nabla^2f_i$，对扰动 KKT 系统作牛顿线性化：
\[
 \begin{bmatrix}
 H&A^{\mathsf T}&J^{\mathsf T}&0\\
 A&0&0&0\\
 J&0&0&I\\
 0&0&S&\Lambda
 \end{bmatrix}
 \begin{bmatrix}\Delta x\\\Delta\nu\\\Delta\lambda\\\Delta s\end{bmatrix}
 =-\begin{bmatrix}r_{\mathrm d}\\r_{\mathrm e}\\r_{\mathrm i}\\r_{\mathrm c}\end{bmatrix}.
\]
在此系统非奇异时求出方向。先取不超过使 $s+t\Delta s>0$、$\lambda+t\Delta\lambda>0$ 的最大步长的固定比例，再回溯保证定义域和残差要求。更新全部变量，重新计算平均互补量，直到可行性残差、驻点残差和互补量均达到要求。

$s^{\mathsf T}\lambda$ 常作为代理对偶间隙。只有原可行性成立且 $x$ 确实最小化对应拉格朗日函数时，它才等于真实的 $f_0(x)-g(\lambda,\nu)$；在不可行迭代点上不能仅凭它很小就宣称近似最优。

保持正变量的步长可以直接计算。令
\[
 t_{\mathrm{pos}}=\min\left\{
 \min_{\Delta s_i<0}\frac{-s_i}{\Delta s_i},
 \ \min_{\Delta\lambda_i<0}\frac{-\lambda_i}{\Delta\lambda_i}
 \right\},
\]
空指标集上的最小值按 $+\infty$ 处理。选取固定 $0<\tau<1$，从 $t=\min\{1,\tau t_{\mathrm{pos}}\}$ 开始回溯，就不会在一步中把正的松弛量或乘子变成零。对线性约束，原残差的线性化是精确的；对非线性凸约束，还须检查更新后的实际约束残差，不能只依赖线性预测。

\section{综合实现与课程总结练习}

把全书应用于带界最小二乘：最小化 $\|Ax-b\|_2^2/2$，满足 $0\leq x\leq u$，其中 $u>0$。这是 QP，$x=u/2$ 为严格可行初值。障碍为 $-\sum_j\log x_j-\sum_j\log(u_j-x_j)$，精确中心点间隙为 $2n/t$；中心化目标的海塞矩阵为 $tA^{\mathsf T}A$ 加正的对角矩阵，因此中心化方向唯一可解。即使 $A$ 秩亏，障碍子问题仍有严格曲率。

记中心化目标为 $F_t(x)=t\|Ax-b\|_2^2/2+\phi(x)$，其梯度与海塞矩阵具体为
\[
 \nabla F_t(x)=tA^{\mathsf T}(Ax-b)
  +\left(-\frac1{x_j}+\frac1{u_j-x_j}\right)_{j=1}^n,
\]
\[
 \nabla^2F_t(x)=tA^{\mathsf T}A+
 \operatorname{diag}\left(\frac1{x_j^2}+\frac1{(u_j-x_j)^2}\right)_{j=1}^n.
\]
初值 $u/2$ 严格位于盒内；每步解海塞矩阵方程得到方向 $d$。保持盒内可行的步长上界为
\[
 \eta_{\max}=\min\left\{
 \min_{d_j<0}\frac{-x_j}{d_j},
 \ \min_{d_j>0}\frac{u_j-x_j}{d_j}
 \right\},
\]
其中 $\eta_{\max}$ 表示线搜索步长的上界，空指标集仍按 $+\infty$ 处理。取 $0<\tau<1$，从步长 $\eta=\min\{1,\tau\eta_{\max}\}$ 开始回溯，再检查中心化目标下降，便得到一套可实际执行的内层迭代。

若 $A$ 满列秩，令 $\mu=\lambda_{\min}(A^{\mathsf T}A)>0$，两组盒约束均仿射，所以原拉格朗日函数的强凸常数仍为 $\mu$。采用当前障碍乘子时，$r=A^{\mathsf T}(Ax-b)+\lambda^{\mathrm{up}}-\lambda^{\mathrm{low}}$，其中 $\lambda_j^{\mathrm{low}}=1/(tx_j)$、$\lambda_j^{\mathrm{up}}=1/(t(u_j-x_j))$。式\eqref{eq:co-inexact-center-bound}于是给出 $2n/t+\|r\|_2^2/(2\mu)$ 的原目标误差上界。若 $A$ 秩亏，内层海塞矩阵仍正定，却不能据此把原拉格朗日函数的强凸常数当作正数；子问题好解与原问题误差可认证是两个条件。

\begin{enumerate}
 \item 独立推导上述模型的中心化梯度、海塞矩阵与步长上界；解释为什么障碍参数增大时仍需重新中心化。
 \item 对 $\min_{0\leq x\leq1}(x-2)^2/2$ 求原问题最优点，写出中心路径方程，解释中心点如何趋于边界。
 \item 举出一个可行但无严格可行点的凸约束系统，说明为什么基本障碍初始化失败不等于不可行。
 \item 完成一个小规模实验：写出模型、凸性依据、KKT 条件和停止标准，记录残差及适用的误差证书。若用软件结果交叉核对，应比较相同目标、约束和精度，而不仅比较变量输出。
\end{enumerate}

\appendix

\chapter{数学基础与符号补充}
\label{app:co-background}

本附录集中列出正文使用的基础工具。首次阅读时可以按引用查阅；各结论的条件也是分析优化模型时应保留的条件。

\section{范数、内积与对偶范数}

范数 $\|x\|$ 满足正定性、绝对齐次性 $\|\alpha x\|=|\alpha|\|x\|$ 和三角不等式。对 $1\leq p<\infty$，$\|x\|_p=(\sum_i|x_i|^p)^{1/p}$；$\|x\|_\infty=\max_i|x_i|$。若 $p,q>1$ 且 $1/p+1/q=1$，Young 不等式 $ab\leq a^p/p+b^q/q$ 对非负 $a,b$ 成立，它可由固定 $a$ 后求右侧减左侧关于 $b$ 的最小值证明。对归一化分量应用并求和，即得 Hölder 不等式
\[
 |x^{\mathsf T}y|\leq\|x\|_p\|y\|_q.
\]
端点形式为 $|x^{\mathsf T}y|\leq\|x\|_1\|y\|_\infty$；$p=q=2$ 为柯西--施瓦茨不等式。

对偶范数 $\|y\|_*=\max_{\|x\|\leq1}y^{\mathsf T}x$ 的最大值在有限维紧单位球上取到。Hölder 不等式给出上界，选取同号且大小成适当幂次的分量可取等号，所以 $\ell_p$ 的对偶为 $\ell_q$。$P\succ0$ 时，$\sqrt{x^{\mathsf T}Px}$ 的对偶范数为 $\sqrt{y^{\mathsf T}P^{-1}y}$。

\section{有限维拓扑与极值存在性}

有限维空间中，集合紧当且仅当闭且有界；紧集中的任意序列都有收敛于该集合的子列。连续函数在非空紧集上取到最大、最小值。较弱的下半连续性要求 $f(x)\leq\liminf_k f(x_k)$ 对所有 $x_k\to x$ 成立；在非空紧集上，它已足以保证有限值函数取到最小值。

\begin{proposition}
设 $f:C\to\mathbb R$ 下半连续，存在实数 $\alpha$ 使次水平集 $S_\alpha=\{x\in C:f(x)\leq\alpha\}$ 非空且紧，则 $f$ 在 $C$ 上有全局最小点。
\end{proposition}
\begin{proof}
先在紧集 $S_\alpha$ 上取极小化序列。若其函数值无下界，紧性给出收敛子列，下半连续性将迫使有限函数值不大于 $-\infty$，矛盾。因此下确界有限，取收敛子列后，下半连续性说明极限点取到该下确界。$S_\alpha$ 外的函数值大于 $\alpha$，故该点也是全局最小点。
\end{proof}

强制性（coercivity）指 $\|x\|_2\to\infty$ 时 $f(x)\to+\infty$。在非空闭集上的连续强制函数有非空紧次水平集，因而取到最小值。开集上的连续函数则可能在缺失的边界逼近下确界，如 $f(x)=x$ 在 $(0,1)$ 上。

\begin{example}[用次水平集证明最优解存在]
设 $C\subseteq\mathbb R^n$ 非空闭，$A,b$ 给定，$\alpha>0$。证明 $f(x)=\|Ax-b\|_2^2/2+\alpha\|x\|_2^2/2$ 在 $C$ 上达到最小值。
\end{example}
\begin{solve}
取任意 $x_0\in C$。在次水平集 $f(x)\leq f(x_0)$ 中，由 $f(x)\geq\alpha\|x\|_2^2/2$ 得 $\|x\|_2\leq\sqrt{2f(x_0)/\alpha}$。该次水平集非空、闭且有界，故紧，连续目标在其上达到最小值。若 $C$ 还凸，强凸性进一步给出唯一性；没有凸性时，上面的存在性证明仍然成立，但唯一性不再有保证。
\end{solve}

\section{微分与矩阵分析}

可微性表示 $f(x+h)=f(x)+\nabla f(x)^{\mathsf T}h+o(\|h\|_2)$，要求余项在所有趋零方向上一致地小。沿每条直线有方向导数并不足够：令 $f(0,0)=0$，非原点处 $f(x,y)=x^3/(x^2+y^2)$，其方向导数 $d_1^3/(d_1^2+d_2^2)$ 存在但不是方向的线性函数，故原点不可微。

二次连续可微时，泰勒展开为 $f(x+h)=f(x)+\nabla f(x)^{\mathsf T}h+\tfrac12h^{\mathsf T}\nabla^2f(x)h+o(\|h\|_2^2)$。仿射复合 $g(z)=f(Az+b)$ 满足 $\nabla g=A^{\mathsf T}\nabla f$、$\nabla^2g=A^{\mathsf T}(\nabla^2f)A$。这些关系解释了降维和变量缩放对曲率的作用。

对称矩阵可谱分解为 $X=U\Lambda U^{\mathsf T}$，$U$ 正交；它半正定当且仅当全部特征值非负。正定时平方根和逆可通过对特征值分别开方、取倒数定义。谱范数为 $\|A\|_2=\max_{\|x\|_2=1}\|Ax\|_2$，等于最大奇异值。正定矩阵的 $\log\det X$ 沿方向 $H$ 的导数为 $\operatorname{tr}(X^{-1}H)$；逆矩阵的方向导数为 $-X^{-1}HX^{-1}$，后者可对恒等式 $XX^{-1}=I$ 求导得到。

\begin{example}[最小二乘梯度的维度与推导]
令 $A\in\mathbb R^{r\times n}$、$x\in\mathbb R^n$、$b\in\mathbb R^r$，求 $f(x)=\|Ax-b\|_2^2/2$ 的梯度和海塞矩阵。
\end{example}
\begin{solve}
设扰动为 $h\in\mathbb R^n$，展开得
\[
 f(x+h)-f(x)=(Ax-b)^{\mathsf T}Ah+\frac12h^{\mathsf T}A^{\mathsf T}Ah.
\]
把一次项写成 $[A^{\mathsf T}(Ax-b)]^{\mathsf T}h$，就读出梯度 $A^{\mathsf T}(Ax-b)\in\mathbb R^n$；二次项给出海塞矩阵 $A^{\mathsf T}A\in\mathbb R^{n\times n}$。残差有 $r$ 个分量，而梯度必须有 $n$ 个分量，维度检查可以立即排除把梯度误写为 $A(Ax-b)$ 的错误。
\end{solve}

实际计算矩阵微分时，先写出扰动的一次项，再按内积读出梯度，通常比逐元素求导更不易出错。正定矩阵的情形也一样：若 $X=U\Lambda U^{\mathsf T}\succ0$，则 $\log\det X=\sum_i\log\lambda_i$，它把各特征方向的尺度相加；当任一特征值趋于零时，$-\log\det X$ 趋于无穷，这解释了它作为半正定锥障碍函数的作用。

\section{概率工具与学习检查}

随机向量 $Z$ 的均值为 $\mu=\mathbb EZ$，协方差为 $\Sigma=\mathbb E[(Z-\mu)(Z-\mu)^{\mathsf T}]$。若二阶矩有限，对给定对称矩阵 $P$ 有
\[
 \mathbb E[Z^{\mathsf T}PZ]=\mu^{\mathsf T}P\mu+\operatorname{tr}(P\Sigma),
\]
可由 $Z=\mu+(Z-\mu)$ 展开并利用中心化项均值为零得到。独立性推出协方差为零，反向一般不成立。

交换参数微分与期望需要条件。例如在参数点的某邻域中，样本函数几乎处处可微，其各偏导绝对值受同一个可积随机变量控制，且函数在某参考参数处可积，则可用支配收敛定理交换微分与期望。计算似然或信息矩阵时，应验证这些条件或使用无需交换的有限和表达。

\begin{enumerate}
 \item 证明 $\|x\|_2\leq\|x\|_1\leq\sqrt n\|x\|_2$，并给出两端取等号的例子。
 \item 构造闭但不有界的可行集及连续目标，使有限下确界不取到。提示：$e^x$ 在实轴上。
 \item 对 $X\succ0$，由逆矩阵导数重新推导 $-\log\det X$ 的二阶方向导数。
\end{enumerate}

\chapter{两个二次函数相关问题}
\label{app:co-quadratic}

本附录为选读。它展示某些非凸二次问题仍有精确的半正定证书；这种特殊结构不能推广为一般非凸优化的结论。

\section{单个二次约束与信赖域问题}

信赖域子问题（trust-region subproblem）为
\[
 \min_{\|x\|_2\leq\Delta}\left\{\frac12x^{\mathsf T}Qx+c^{\mathsf T}x\right\},
 \qquad Q=Q^{\mathsf T},\quad\Delta>0.
\]
即使 $Q$ 不半正定，连续目标在紧球上仍取到最小值。其全局最优点可由以下条件刻画：存在 $\lambda\geq0$，使
\[
 (Q+\lambda I)x=-c,\quad Q+\lambda I\succeq0,\quad
 \|x\|_2\leq\Delta,\quad\lambda(\|x\|_2^2-\Delta^2)=0.
\]
充分性可直接验证：拉格朗日函数 $\mathcal L(x,\lambda)=x^{\mathsf T}Qx/2+c^{\mathsf T}x+\lambda(\|x\|_2^2-\Delta^2)/2$ 在该点全局最小，且互补性使其值等于原目标，弱对偶给出最优性。必要性由下一节的 S 引理推出，这里引用该标准结论。

\begin{example}
令 $Q=\operatorname{diag}(-2,2)$、$c=0$、$\Delta=1$，目标为 $-x_1^2+x_2^2$。其最小值为 $-1$，在 $(1,0)$、$(-1,0)$ 取到；取 $\lambda=2$ 时 $Q+\lambda I=\operatorname{diag}(0,4)\succeq0$，符合上述条件。这也说明最优点不必唯一，修正海塞矩阵也不必正定。
\end{example}

\section{S 引理与半正定证书}

\begin{theorem}[S 引理]
设 $q_0,q_1$ 为实二次函数，存在 $\bar x$ 使 $q_1(\bar x)<0$。则下列两项等价：对所有 $x$，$q_1(x)\leq0$ 蕴含 $q_0(x)\geq0$；存在 $\lambda\geq0$，使 $q_0(x)+\lambda q_1(x)\geq0$ 对所有 $x$ 成立。
\end{theorem}
由第二项推出第一项只需移项；反向是引理的实质，本书引用其结论，证明可参见\cite{ahmadi-slemma}。严格可行性不可直接删除，多约束时通常也只保留乘子证书的充分性。

写 $q_i(x)=x^{\mathsf T}Q_i x+2a_i^{\mathsf T}x+c_i$，其全空间非负性证书等价于
\[
 \begin{bmatrix}
 Q_0+\lambda Q_1&a_0+\lambda a_1\\
 (a_0+\lambda a_1)^{\mathsf T}&c_0+\lambda c_1
 \end{bmatrix}\succeq0.
\]
这把无限多个点上的不等式化成有限维矩阵条件。二次函数全空间非负等价于相应提升矩阵半正定：末坐标非零的向量可缩放为 $(x,1)$，末坐标为零的情形由连续性补上。将 $q_0$ 换为 $q_0-\gamma$，最大化可认证的下界 $\gamma$，就得到单约束二次问题的精确 SDP 值表示。

在信赖域问题中，可以把这一必要性论证写完整。令原目标为 $f$、最优值为 $p^\star$，取 $q_0=f-p^\star$、$q_1=(\|x\|_2^2-\Delta^2)/2$。原点严格满足 $q_1<0$，且在球上 $q_0\geq0$，S 引理给出 $\lambda\geq0$ 使 $q_0+\lambda q_1$ 在全空间非负。在任意原最优点 $x^\star$，这个非负量又等于 $\lambda q_1(x^\star)\leq0$，故它为零并满足互补松弛。非负二次函数在这个零点取全局最小值，因此其梯度为零、二次项半正定，恰好得到前述完整条件。这个推导明确说明了单约束的严格可行性在哪一步发挥作用。

\section{联合值域与应用练习}

对两个对称矩阵 $A,B$，齐次二次型的联合值域 $\{(x^{\mathsf T}Ax,x^{\mathsf T}Bx):x\in\mathbb R^n\}$ 是凸锥；这是两二次型的经典凸性结论，此处引用，详见\cite{boyd-vandenberghe}附录中的论述。该锥未必闭，例如 $(x^2,2xy)$ 的值域包含所有第一坐标正的点及原点，却不包含 $(0,1)$。

不能将此结论直接推广到任意非齐次二次函数：$\{(x,x^2):x\in\mathbb R\}$ 是抛物线，两个不同点的中点通常不在其上。单约束 S 引理中的资格条件正是把联合值域几何与分离证书联系起来所需的关键假设之一。

\begin{enumerate}
 \item 用 $q_1(x)=x^2-1$、$q_0(x)=1-x^2$ 验证 S 引理，给出乘子 $\lambda=1$。
 \item 取 $q_1(x)=x^2$、$q_0(x)=x$。验证蕴含关系成立，却不存在有限 $\lambda\geq0$ 使 $x+\lambda x^2$ 对所有实数非负，说明严格可行性的作用。
 \item 对球内二次函数的下界认证写出 SDP，并与直接计算的最小特征值例题核对。
\end{enumerate}

\chapter{数值线性代数与算法实现基础}
\label{app:co-numerical}

优化算法经常把主要成本集中在线性方程求解。数值实现应利用正定性、稀疏性和秩结构，同时区分方程残差与实际解误差。

\section{线性系统、分解与最小二乘}

对可逆方阵，带主元 LU 分解将置换后的矩阵写成下三角与上三角矩阵之积，再作两次三角求解。对称正定矩阵可写成 $LL^{\mathsf T}$，其中 $L$ 为对角元正的下三角矩阵，称为 Cholesky 分解。对称不定矩阵通常用带置换的 $LDL^{\mathsf T}$ 分解，$D$ 可含 $1\times1$ 与 $2\times2$ 对角块。KKT 矩阵即使可逆，也通常不正定，因此不能直接套用普通 Cholesky 分解。

对 $A\in\mathbb R^{r\times n}$ 满列秩的最小二乘，薄 QR 分解为 $A=QR$，$Q^{\mathsf T}Q=I$、$R$ 为可逆上三角矩阵。分解残差可得
\[
 \|Ax-b\|_2^2=\|Rx-Q^{\mathsf T}b\|_2^2+\|(I-QQ^{\mathsf T})b\|_2^2,
\]
所以只需解 $Rx=Q^{\mathsf T}b$。

一般矩阵的奇异值分解为 $A=U\Sigma V^{\mathsf T}$。将非零奇异值取倒数并转置矩形对角结构，得到 $\Sigma^\dagger$，定义 Moore--Penrose 伪逆 $A^\dagger=V\Sigma^\dagger U^{\mathsf T}$。$A^\dagger b$ 是所有最小二乘解中的最小范数解；若方程可行，它就是最小范数精确解。

\begin{example}[秩亏时拟合值可以唯一而参数不唯一]
令 $A=\left[\begin{smallmatrix}1&1\\2&2\end{smallmatrix}\right]$、$b=(1,2)^{\mathsf T}$。所有满足 $x_1+x_2=1$ 的点都使残差为零，因此参数并不唯一，拟合向量 $Ax=b$ 却唯一。把任意解写成 $(1/2,1/2)+(d,-d)$，其平方范数为 $1/2+2d^2$，所以伪逆选择的最小范数解为 $(1/2,1/2)$。不能因为算法返回了这个特定向量，就认为数据唯一确定了两个参数。
\end{example}

\section{块矩阵、Schur 补与 KKT 方程}

对块方程 $Hu+A^{\mathsf T}v=r$、$Au=s$，若 $H$ 可逆，先写 $u=H^{-1}(r-A^{\mathsf T}v)$，再得到 $AH^{-1}A^{\mathsf T}v=AH^{-1}r-s$。式中负号必须与原方程一致；优化中的右端常由负残差组成，机械套公式容易出错。

若 $H\succ0$ 且 $A$ 满行秩，Schur 补正定；若 $H$ 奇异但在 $\ker A$ 上正定，整个 KKT 矩阵仍可能可逆，此时直接的块求逆公式不可用。稀疏直接分解中的填充是原零元素在分解中变为非零的现象，适当排序可降低存储和运算成本。

\section{条件数、舍入误差与规模}

对可逆矩阵 $A$，条件数 $\kappa_2(A)=\|A\|_2\|A^{-1}\|_2$ 衡量输入扰动被放大的程度。若 $Ax=b$、近似解为 $\widehat x$、残差为 $r=b-A\widehat x$，则
\[
 \|x-\widehat x\|_2\leq\|A^{-1}\|_2\|r\|_2,
 \qquad
 \frac{\|x-\widehat x\|_2}{\|x\|_2}
 \leq\kappa_2(A)\frac{\|r\|_2}{\|b\|_2}
\]
在 $x,b\ne0$ 时成立。因此小相对残差未必意味着小相对解误差。后向误差讨论计算解可看成哪个邻近问题的精确解；条件数则决定邻近问题的解与原解能差多少。

例如 $A=\operatorname{diag}(1,10^{-8})$、$b=(1,10^{-8})^{\mathsf T}$ 的真解为 $(1,1)$。近似解 $\widehat x=(1,0)$ 的残差仅为 $(0,10^{-8})^{\mathsf T}$，但第二个分量的误差为一。原因是第二个方向的系数很小，变量的大变化只引起方程左侧的小变化。这与平坦优化目标中小梯度不保证点误差小，是同一种尺度问题的不同表现。

满列秩矩阵满足 $\kappa_2(A^{\mathsf T}A)=\kappa_2(A)^2$。正规方程虽然形式简洁，却可能放大数值困难；QR 和奇异值分解避免显式形成这一平方条件的矩阵。不同方法的实际误差仍与舍入、缩放及算法实现有关。

稠密方阵分解一般需 $O(n^3)$ 运算，已有分解后的每个新右端求解需 $O(n^2)$。稀疏问题取决于非零结构与填充，不能只用原非零元素个数给出统一成本。牛顿法和内点法中若矩阵结构不变，可以复用符号分解；数值因子是否可复用则取决于矩阵值是否改变。

\section{综合练习与结果核验}
\begin{enumerate}
 \item 对 $A=\operatorname{diag}(1,10^{-8})$，构造残差很小但解误差明显的近似解，核对条件数界。
 \item 对同一个满列秩最小二乘问题，比较正规方程、QR 与奇异值分解的结果；逐步使两列接近线性相关，记录残差与参数误差的差别。
 \item 任选一次牛顿或 KKT 步，核验方程残差、方向下降性、更新后可行性及优化停止指标。线性系统求解成功、迭代收敛和达到所需最优性精度是三个需要分别确认的结论。
\end{enumerate}

\begin{thebibliography}{9}
\bibitem{boyd-vandenberghe}
Stephen Boyd and Lieven Vandenberghe.
\emph{Convex Optimization}. Cambridge University Press, 2004.
作者提供的教材、勘误与教学资料：\url{https://web.stanford.edu/~boyd/cvxbook/}。

\bibitem{parikh-boyd}
Neal Parikh and Stephen Boyd.
\emph{Proximal Algorithms}.
Foundations and Trends in Optimization, 1(3):123--231, 2014.
作者提供的论文与示例：\url{https://web.stanford.edu/~boyd/papers/prox_algs.html}。


\bibitem{ahmadi-slemma}
Amir Ali Ahmadi; Georgina Hall（课堂记录）.
\emph{ORF 523: Convex and Conic Optimization, Lecture 12}.
Princeton University, Spring 2017.
\url{https://aaa.princeton.edu/orf523s17}。
\end{thebibliography}


\end{document}
